\begin{derivation}[颜色平均、自旋迹与两次 Jacobian]
虚光子不改变颜色。取夸克和反夸克颜色为 $i,j$，其振幅含 $\delta_{ij}$，故
\begin{equation*}
 \frac1{N_c^2}\sum_{i,j}|\delta_{ij}|^2=\frac1{N_c}.
\end{equation*}
它与 $e^+e^-\to q\bar q$ 的末态颜色求和 $N_c$ 不同；交叉振幅并不会自动把初态平均一起交叉。
定义
\begin{equation*}
 T^{\mu\nu}(a,b)=\Tr(\slashed a\gamma^\mu\slashed b\gamma^\nu)
 =4(a^\mu b^\nu+a^\nu b^\mu-\eta^{\mu\nu}a\cdot b).
\end{equation*}
两迹的指标缩并为
\begin{equation*}
 T^{\mu\nu}(p_2,p_1)T_{\mu\nu}(k_-,k_+)
 =32[(p_2\cdot k_-)(p_1\cdot k_+)+(p_2\cdot k_+)(p_1\cdot k_-)]
 =8(\hat u^2+\hat t^2).
\end{equation*}
再乘初态自旋平均 $1/4$、颜色平均 $1/N_c$ 以及 $g_{\rm em}^4Q_f^2/\hat s^2$，便得振幅平方。用 $z_\theta=\cos\theta$，
\begin{align*}
 \hat t&=-\frac{\hat s}{2}(1-z_\theta),\qquad
 \hat u=-\frac{\hat s}{2}(1+z_\theta),\\
 \frac{\dd\hat\sigma}{\dd z_\theta}
 &=\frac{\hbar^2}{32\pi\hat s}\overline{|\mathcal M|^2}
 =\frac{\pi\alpha^2\hbar^2Q_f^2}{2N_c\hat s}(1+z_\theta^2).
\end{align*}
由于 $\int_{-1}^{1}\dd z_\theta(1+z_\theta^2)=8/3$，角积分给\eqnrefs{eq:dy-born-partonic}。

将 Born 截面乘上两条运动学分布
\begin{equation*}
 \delta(\mathcal Q^2-sx_1x_2)
 \delta\!\left(Y-\frac12\ln\frac{x_1}{x_2}\right).
\end{equation*}
积分时需要的行列式为
\begin{equation*}
 \det\frac{\partial(sx_1x_2,\frac12\ln(x_1/x_2))}{\partial(x_1,x_2)}
 =\det\begin{pmatrix}sx_2&sx_1\\1/(2x_1)&-1/(2x_2)\end{pmatrix}=-s.
\end{equation*}
因此两次 delta 积分产生 $1/s$，而不是 $1/(x_1x_2s)$。部分子截面中的 $1/\hat s$ 则在支撑点变为 $1/\mathcal Q^2$，二者合并给\eqnrefs{eq:dy-born-hadronic}。其量纲为面积除以动量平方，与左边完全一致。
\end{derivation}

\begin{derivation}[从\eqnrefs{eq:dis-W-decomposition-dp13}到\eqnrefs{eq:dis-tree-parton-f2-dp76}]
在领先扭曲、无质量部分子极限中，被击中夸克的入射动量写成
\begin{equation*}
 p^\mu=xP^\mu,
 \qquad p^2=0,
\end{equation*}
而吸收虚光子后的末态夸克满足 $(p+q)^2=0$。因此
\begin{equation*}
 0=(p+q)^2=2p\cdot q-\mathcal Q^2
 \quad\Longrightarrow\quad
 p\cdot q=\frac{\mathcal Q^2}{2}.
\end{equation*}
另一方面 $p=xP$，所以
\begin{equation*}
 x\,P\cdot q=\frac{\mathcal Q^2}{2}
 \quad\Longrightarrow\quad x=x_B.
\end{equation*}
更严格地，相空间中的质量壳分布满足
\begin{equation*}
 \delta\!\bigl((xP+q)^2\bigr)
 =\delta\!\bigl(2P\cdot q\,(x-x_B)\bigr)
 =\frac{x_B}{\mathcal Q^2}\,\delta(x-x_B),
\end{equation*}
其中已忽略 $x^2M_H^2c^2/\mathcal Q^2$ 的靶质量修正。

树级电磁流的自旋平均张量分子为
\begin{align*}
 t_{\mu\nu}(p,q)
 &\equiv \frac12\Tr\!\left[\slashed p\gamma_\mu(\slashed p+\slashed q)\gamma_\nu\right]\\
 &=2\left[p_\mu(p+q)_\nu+p_\nu(p+q)_\mu
 -\eta_{\mu\nu}p\cdot(p+q)\right].
\end{align*}
由 $p^2=0$ 与 $p\cdot q=\mathcal Q^2/2$，有 $p\cdot(p+q)=\mathcal Q^2/2$。定义与 $q^\mu$ 横向的部分子动量
\begin{equation*}
 \widetilde p_\mu
 \equiv p_\mu-\frac{p\cdot q}{q^2}q_\mu
 =p_\mu+\frac12q_\mu,
\end{equation*}
于是 $p_\mu=\widetilde p_\mu-q_\mu/2$、$(p+q)_\mu=\widetilde p_\mu+q_\mu/2$。代入上式逐项整理：
\begin{align*}
 p_\mu(p+q)_\nu+p_\nu(p+q)_\mu
 &=2\widetilde p_\mu\widetilde p_\nu-\frac12q_\mu q_\nu,\\
 t_{\mu\nu}
 &=4\widetilde p_\mu\widetilde p_\nu-q_\mu q_\nu-\mathcal Q^2\eta_{\mu\nu}\\
 &=\mathcal Q^2\Pi_{\mu\nu}^{(q)}+4\widetilde p_\mu\widetilde p_\nu.
\end{align*}
最后一步使用 $q^2=-\mathcal Q^2$ 与正文\eqnrefs{eq:dis-transverse-tensor-basis-dp68}。因为 $p=x_BP$，横向投影也满足 $\widetilde p_\mu=x_B\widetilde P_\mu$，故单个味的张量必为
\begin{equation*}
 w_{\mu\nu}^{(f)}
 =\mathcal C_f(x_B)
 \left[
 \mathcal Q^2\Pi_{\mu\nu}^{(q)}
 +4x_B^2\widetilde P_\mu\widetilde P_\nu
 \right].
\end{equation*}
比较正文\eqnrefs{eq:dis-W-decomposition-dp13} 的两个独立张量系数，立即有
\begin{equation*}
 F_1^{(f)}=\mathcal C_f\mathcal Q^2,
 \qquad
 \frac{F_2^{(f)}}{P\cdot q}=4x_B^2\mathcal C_f.
\end{equation*}
再用 $P\cdot q=\mathcal Q^2/(2x_B)$ 消去 $\mathcal C_f$：
\begin{equation*}
 \boxed{F_2^{(f)}=2x_BF_1^{(f)}}.
\end{equation*}
因此 Callan--Gross 关系由无质量自旋 $1/2$ 电磁流的 Dirac 迹与质量壳条件共同确定。

还需固定整体归一化。按正文\eqnrefs{eq:pdf-operator-definition-dp14} 的 PDF 归一化，树级系数函数对应
\begin{equation*}
 F_1^{(f)}(x_B)=\frac{Q_f^2}{2}
 \left[f_{f/H}(x_B,\mu_{\rm F})+f_{\bar f/H}(x_B,\mu_{\rm F})\right].
\end{equation*}
将它代入刚得到的 Callan--Gross 关系，便有
\begin{equation*}
 F_2^{(f)}(x_B)
 =x_BQ_f^2
 \left[f_{f/H}(x_B,\mu_{\rm F})+f_{\bar f/H}(x_B,\mu_{\rm F})\right].
\end{equation*}
最后对所有活跃味求和，得到正文\eqnrefs{eq:dis-tree-parton-f2-dp76}，同时正文\eqnrefs{eq:dis-longitudinal-structure-function-dp31} 给出 $F_L=F_2-2x_BF_1=0$。非零 $F_L$ 因而必须来自夸克质量、横向 QCD 辐射或更高扭曲度等超出这一极限的效应。
\end{derivation}

\begin{derivation}[从\eqnrefs{eq:dis-lightlike-wilson-line-dp68}到\eqnrefs{eq:pdf-operator-definition-dp14}]
部分子模型中的 $f_{f/H}$ 必须提升为 QCD 中的规范不变算符矩阵元。取一个指向未来的无量纲类光向量 $n^\mu$，满足

\begin{equation*}
 n^2=0,
 \qquad P\cdot n>0.
\end{equation*}
沿 $n^\mu$ 的实际时空间隔写成
\begin{equation*}
 z^\mu(\tau)
 =\frac{\hbar\tau}{P\cdot n}\,n^\mu,
\end{equation*}
其中 $\tau$ 无量纲；$\hbar/(P\cdot n)$ 具有长度量纲，因此 $z^\mu$ 保持 SI 长度。

两个不同点的夸克双线性本身不规范不变。代入正文\eqnrefs{eq:dis-lightlike-wilson-line-dp68} 的基本表示 Wilson 线
\begin{equation*}
 \mathcal W_n(\tau,0)
 \equiv\mathcal P\exp\left[
 \ii g_s\int_0^{\hbar\tau/(P\cdot n)}\dd\lambda\,
 n^\mu\mathscr C_\mu^a(\lambda n)T^a
 \right].
\end{equation*}
因为 $[\mathscr C_\mu^a]=\mathrm{m^{-1}}$、$[\dd\lambda]=\mathrm m$，指数严格无量纲。在局域颜色转动 $U(x)$ 下，规范联络的变换使路径有序指数满足
\begin{equation*}
 \mathcal W_n(\tau,0)
 \longrightarrow
 U[z(\tau)]\mathcal W_n(\tau,0)U^\dagger(0),
\end{equation*}
而夸克场满足
\begin{equation*}
 \psi_f[z(\tau)]\to U[z(\tau)]\psi_f[z(\tau)],
 \qquad
 \bar\psi_f[z(\tau)]\to\bar\psi_f[z(\tau)]U^\dagger[z(\tau)].
\end{equation*}
因此
\begin{equation*}
 \bar\psi_f[z(\tau)]\slashed n\,
 \mathcal W_n(\tau,0)\psi_f(0)
\end{equation*}
是颜色单态。

在协变强子态归一化下，把分子矩阵元与 $2P\cdot n$ 的比值暂时记作
\begin{equation*}
 \mathcal B_{f/H}(\tau,\mu_{\rm F})
 \equiv
 \frac{\langle P|
 \bar\psi_f[z(\tau)]\slashed n\mathcal W_n(\tau,0)\psi_f(0)
 |P\rangle_{\mu_{\rm F}}}
 {2P\cdot n}.
\end{equation*}
其无量纲 Fourier 变换就是夸克部分子分布
\begin{equation*}
 f_{f/H}(x,\mu_{\rm F})
 =\int_{-\infty}^{+\infty}
 \frac{\dd\tau}{2\pi}\,
 \ee^{-\ii x\tau}
 \mathcal B_{f/H}(\tau,\mu_{\rm F}).
\end{equation*}
\end{derivation}

\begin{derivation}[从\eqnrefs{eq:dis-current-product-ope-dp68}到\eqnrefs{eq:dis-ope-moment-factorization-dp13}]
深度非弹性极限由 $\mathcal Q^2\gg\Lambda_{\rm had}^2$ 控制；Fourier 积分的主贡献来自 $|z^2|=O(\hbar^2/\mathcal Q^2)$ 的短距离区域。对任意选定的有限 OPE 截断集合 $\mathscr S$，先把余项定义为
\begin{equation*}
 \mathcal R_{\mu\nu}^{(\mathscr S)}(z)
 \equiv
 \mathcal T\{j_\mu(z)j_\nu(0)\}
 -\sum_{(N,i)\in\mathscr S}
 C_{\mu\nu,N}^{(i)}(z,\mu_{\rm F})
 \mathcal O_N^{(i)}(0,\mu_{\rm F}).
\end{equation*}
于是有限阶关系是严格等式
\begin{equation*}
 \mathcal T\{j_\mu(z)j_\nu(0)\}
 =
 \sum_{(N,i)\in\mathscr S}
 C_{\mu\nu,N}^{(i)}(z,\mu_{\rm F})
 \mathcal O_N^{(i)}(0,\mu_{\rm F})
 +\mathcal R_{\mu\nu}^{(\mathscr S)}(z).
\end{equation*}
OPE 的动力学内容是：按扭曲度与算符维数扩大 $\mathscr S$ 时，$\mathcal R_{\mu\nu}^{(\mathscr S)}$ 在上述短距离区域按相应的 $\Lambda_{\rm had}^2/\mathcal Q^2$ 幂次受抑；因此截断误差与所保留的 twist 层级一一对应。
领先的非极化夸克算符是对称无迹、扭曲度为 $2$ 的张量
\begin{equation*}
 \mathcal O_{f,N}^{\mu_1\cdots\mu_N}
 =\bar\psi_f\gamma^{\{\mu_1}
 (\ii\hbar D^{\mu_2})\cdots
 (\ii\hbar D^{\mu_N\}})\psi_f
 -\text{迹项}.
\end{equation*}
每个 $\ii\hbar D^\mu$ 都具有物理动量量纲。Lorentz 协变性把前向矩阵元压缩为一个约化标量系数
\begin{equation*}
 \langle P|\mathcal O_{f,N}^{\mu_1\cdots\mu_N}|P\rangle
 =2A_{f,N}(\mu_{\rm F})
 P^{\{\mu_1}\cdots P^{\mu_N\}}
 -\text{迹项},
\end{equation*}
按此定义 $A_{f,N}$ 无量纲。

PDF 与这些局域算符之间的联系来自类光算符对 $\tau$ 的规范协变 Taylor 展开。先把 Wilson 线与被平移的场搬到同一端点：
\begin{equation*}
 \mathcal W_n(0,\tau)\psi_f[z(\tau)]
 =\sum_{r=0}^{\infty}
 \frac1{r!}
 \left(\frac{\hbar\tau}{P\cdot n}\right)^r
 (n\cdot D)^r\psi_f(0).
\end{equation*}
对厄米共轭双局域算符同样展开，便逐阶产生带协变导数的局域算符。另一方面
\begin{equation*}
 x^{N-1}\ee^{-\ii x\tau}
 =(\ii\partial_\tau)^{N-1}\ee^{-\ii x\tau}.
\end{equation*}
把 $\partial_\tau$ 通过分部积分移到 $\mathcal B_{f/H}$ 上，每次导数都带出 $\ii\hbar n\cdot D/(P\cdot n)$。因此 Mellin 矩与局域矩阵元一一对应：
\begin{equation*}
 \int_0^1\dd x\,x^{N-1}
 \left[f_{f/H}(x,\mu_{\rm F})+(-1)^Nf_{\bar f/H}(x,\mu_{\rm F})\right]
 =A_{f,N}(\mu_{\rm F}).
\end{equation*}
所以局域 OPE 的算符混合与 DGLAP 的 PDF 演化并不是两套独立理论；它们分别是同一类光算符重整化在 Mellin 空间与 $x$ 空间中的表示。
\end{derivation}

\begin{derivation}[从\eqnrefs{eq:dis-ope-moment-factorization-dp13}到\eqnrefs{eq:dis-collinear-factorization-dp14}]
定义 Mellin 卷积

\begin{equation*}
 [A\otimes B](x)
 \equiv\int_x^1\frac{\dd z}{z}
 A(z)B\!\left(\frac{x}{z}\right).
\end{equation*}
它满足结合律：
\begin{align*}
 [A\otimes(B\otimes C)](x)
 &=\int_x^1\frac{\dd z}{z}A(z)
 \int_{x/z}^1\frac{\dd u}{u}B(u)
 C\!\left(\frac{x}{zu}\right),\\
 &=[(A\otimes B)\otimes C](x),
\end{align*}
第二行只交换积分次序并令乘积变量 $v=zu$。这个代数性质保证因子化反项可以在系数函数与 PDF 之间一致移动。

在 $d=4-2\epsilon$ 维正则化中，裸类光算符发生混合，写成卷积矩阵
\begin{equation*}
 f_i^{\rm bare}
 =\sum_j Z_{ij}^{\rm F}\otimes f_j(\mu_{\rm F}).
\end{equation*}
在 $\overline{\rm MS}$ 方案的领先阶，
\begin{equation*}
 Z_{ij}^{\rm F}(z)
 =\delta_{ij}\delta(1-z)
 +\frac{\alpha_s}{2\pi}
 \frac1{\bar\epsilon}
 P_{ij}^{(0)}(z)
 +\order(\alpha_s^2),
\end{equation*}
其中 $1/\bar\epsilon=1/\epsilon-\gamma_E+\ln4\pi$。部分子级系数函数中的相反共线极点由 $(Z^{\rm F})^{-1}$ 消去，留下有限的短距离系数 $C_{a,i}$。于是领先幂结构函数成为
\begin{equation*}
 \mathfrak F_a(x_B,\mathcal Q^2)
 =\sum_i
 \left[
 C_{a,i}\!\left(
 \frac{\mathcal Q^2}{\mu_{\rm F}^2},
 \frac{\mu_{\rm DR}^2}{\mu_{\rm F}^2},
 \alpha_s(\mu_{\rm DR})
 \right)
 \otimes f_{i/H}(\mu_{\rm F})
 \right](x_B)
 +\order\!\left(\frac{\Lambda_{\rm QCD}^2}{\mathcal Q^2}\right).
\end{equation*}
这就是部分子模型在 QCD 中的受控推广：$C_{a,i}$ 可由微扰理论计算，$f_{i/H}$ 是规范不变的非微扰强子矩阵元，而两者对人为因子化尺度的依赖必须互相抵消。
\end{derivation}

\begin{derivation}[从\eqnrefs{eq:pdf-operator-definition-dp14}到\eqnrefs{eq:dglap-general-evolution-dp68}]
PDF 是类光算符的强子矩阵元，因此其紫外/共线重整化不是普通数乘，而是在纵向动量分数上的卷积。写成
\begin{equation*}
 f_i^{\rm bare}=Z^{\rm F}_{ij}\otimes f_j,
 \qquad
 [A\otimes B](x)=\int_x^1\frac{\dd z}{z}A(z)B(x/z).
\end{equation*}
这里矩阵乘法同时包含部分子种类求和与 $x$ 卷积。卷积单位元是
$\delta_{ij}\delta(1-z)$。裸算符不依赖人为选择的因子化尺度，所以
\begin{align*}
 0
 &=\frac{\dd f_i^{\rm bare}}{\dd\ln\mu_{\rm F}^2}\\
 &=\frac{\dd Z^{\rm F}_{ij}}{\dd\ln\mu_{\rm F}^2}\otimes f_j
 +Z^{\rm F}_{ij}\otimes
 \frac{\partial f_j}{\partial\ln\mu_{\rm F}^2}.
\end{align*}
左卷积乘 $(Z^{\rm F})^{-1}$ 后得到精确重整化群形式
\begin{equation*}
 \frac{\partial f_i}{\partial\ln\mu_{\rm F}^2}
 =\gamma^{\rm F}_{ij}\otimes f_j,
 \qquad
 \gamma^{\rm F}
 \equiv-(Z^{\rm F})^{-1}\otimes
 \frac{\dd Z^{\rm F}}{\dd\ln\mu_{\rm F}^2}.
\end{equation*}

关键一步是说明为什么 $Z^{\rm F}$ 中的 $1/\epsilon$ 极点在四维极限后会留下有限导数。维数正则化取 $d=4-2\epsilon$。在最低阶最小减除型方案中可写
\begin{equation*}
 Z^{\rm F}_{ij}(z)
 =\delta_{ij}\delta(1-z)
 +\frac{\alpha_s}{2\pi\epsilon}P_{ij}^{(0)}(z)
 +\order(\alpha_s^2),
\end{equation*}
其中 $P_{ij}^{(0)}$ 就是类光算符一圈共线极点的残数。由于裸耦合含 $(\mu_{\rm F}^2)^\epsilon$，保持裸耦合不变时
\begin{equation*}
 \frac{\dd\alpha_s}{\dd\ln\mu_{\rm F}^2}
 =-\epsilon\alpha_s+\order(\alpha_s^2).
\end{equation*}
因此
\begin{align*}
 \frac{\dd Z^{\rm F}}{\dd\ln\mu_{\rm F}^2}
 &=\frac{1}{2\pi\epsilon}
 \frac{\dd\alpha_s}{\dd\ln\mu_{\rm F}^2}P^{(0)}
 +\order(\alpha_s^2)\\
 &=-\frac{\alpha_s}{2\pi}P^{(0)}
 +\order(\alpha_s^2,\epsilon).
\end{align*}
同时
$(Z^{\rm F})^{-1}=I+\order(\alpha_s/\epsilon)$；在求一圈有限反常维数时，逆矩阵的非平凡部分只贡献更高阶项。于是取 $\epsilon\to0$ 后
\begin{equation*}
 \gamma^{\rm F}_{ij}(z)
 =\frac{\alpha_s(\mu_{\rm F})}{2\pi}P_{ij}^{(0)}(z)
 +\order(\alpha_s^2),
\end{equation*}
从而得到正文 DGLAP 方程
\begin{equation*}
 \frac{\partial f_i(x,\mu_{\rm F}^2)}{\partial\ln\mu_{\rm F}^2}
 =\frac{\alpha_s(\mu_{\rm F})}{2\pi}
 \sum_j[P_{ij}^{(0)}\otimes f_j](x)
 +\order(\alpha_s^2).
\end{equation*}

最后检查人为尺度确实从可观测量中消失。若
$\mathfrak F_a=\sum_i C_{a,i}\otimes f_i$，则
\begin{align*}
 0=\frac{\dd F_a}{\dd\ln\mu_{\rm F}^2}
 =\sum_i\left[
 \frac{\partial C_{a,i}}{\partial\ln\mu_{\rm F}^2}\otimes f_i
 +C_{a,i}\otimes\gamma^{\rm F}_{ij}\otimes f_j
 \right].
\end{align*}
对任意 PDF 都成立要求
\begin{equation*}
 \frac{\partial C_{a,i}}{\partial\ln\mu_{\rm F}^2}
 =-\sum_j C_{a,j}\otimes\gamma^{\rm F}_{ji}.
\end{equation*}
因此 DGLAP 的物理含义不是“PDF 随能量真的被人为改写”，而是当短程/共线的分界尺度移动时，长程 PDF 与短程系数函数以相反方向重分配同一物理内容；可观测结构函数保持不变。
\end{derivation}



\begin{derivation}[从\eqnrefs{eq:dglap-general-evolution-dp68}到\eqnrefs{eq:dis-f2-scaling-violation-dp32}]
对每个活跃夸克味 $f$，领先阶 DGLAP 方程分别为
\begin{align*}
 \frac{\partial q_f}{\partial\ln\mu^2}
 &=\frac{\alpha_s}{2\pi}
 \left(P_{qq}\otimes q_f+P_{qg}\otimes g\right),\\
 \frac{\partial \bar q_f}{\partial\ln\mu^2}
 &=\frac{\alpha_s}{2\pi}
 \left(P_{qq}\otimes \bar q_f+P_{qg}\otimes g\right),
\end{align*}
而胶子分布满足
\begin{equation*}
 \frac{\partial g}{\partial\ln\mu^2}
 =\frac{\alpha_s}{2\pi}
 \left[
 P_{gg}\otimes g
 +P_{gq}\otimes\sum_{f=1}^{n_f}(q_f+\bar q_f)
 \right].
\end{equation*}
定义味单态组合
\begin{equation*}
 \Sigma(x,\mu)\equiv
 \sum_{f=1}^{n_f}[q_f(x,\mu)+\bar q_f(x,\mu)].
\end{equation*}
把前两式逐味相加时，每个味有一个 $q_f$ 和一个 $\bar q_f$ 接收 $g\to q_f\bar q_f$ 的馈入，因此
\begin{align*}
 \frac{\partial\Sigma}{\partial\ln\mu^2}
 &=\sum_f\left(
 \frac{\partial q_f}{\partial\ln\mu^2}
 +\frac{\partial\bar q_f}{\partial\ln\mu^2}
 \right)\\
 &=\frac{\alpha_s}{2\pi}
 \left[P_{qq}\otimes\Sigma
 +2n_fP_{qg}\otimes g\right].
\end{align*}
与胶子方程并排写成
\begin{equation*}
 \frac{\partial}{\partial\ln\mu^2}
 \binom{\Sigma}{g}
 =\frac{\alpha_s}{2\pi}
 \begin{pmatrix}
 P_{qq}&2n_fP_{qg}\\
 P_{gq}&P_{gg}
 \end{pmatrix}\otimes
 \binom{\Sigma}{g},
\end{equation*}
即正文\eqnrefs{eq:dis-singlet-gluon-dglap-dp31}。这个矩阵就是一般 DGLAP 方程变换到味单态--胶子基后的演化核。作为检查，任取非单态差
\begin{equation*}
 q_i+\bar q_i-q_j-\bar q_j,
\end{equation*}
两份相同的 $P_{qg}\otimes g$ 严格相消，所以非单态扇区闭合，而单态必然与胶子混合。

现在把同一演化投影到可测的电磁结构函数。取 $\mu_{\rm F}=\mathcal Q$，树级硬系数给
\begin{equation*}
 F_2(x,\mathcal Q^2)
 =x\sum_f Q_f^2[q_f(x,\mathcal Q)+\bar q_f(x,\mathcal Q)]
 +\order(\alpha_s).
\end{equation*}
对 $\ln\mathcal Q^2$ 求导时，树级系数本身没有显式尺度导数，因此到 $O(\alpha_s)$ 只需作用在 PDF 上：
\begin{align*}
 \frac{\partial F_2}{\partial\ln\mathcal Q^2}
 &=x\sum_fQ_f^2
 \left(
 \frac{\partial q_f}{\partial\ln\mathcal Q^2}
 +\frac{\partial\bar q_f}{\partial\ln\mathcal Q^2}
 \right)+\order(\alpha_s^2)\\
 &=\frac{\alpha_s}{2\pi}x\sum_fQ_f^2
 \left[
 P_{qq}\otimes(q_f+\bar q_f)
 +2P_{qg}\otimes g
 \right]+\order(\alpha_s^2),
\end{align*}
即正文\eqnrefs{eq:dis-f2-scaling-violation-dp32}。这里的胶子项不是因为电磁流直接耦合胶子，而是因为因子化尺度变化会通过 $g\to q\bar q$ 把胶子分布旋转进带电夸克算符；这正是算符混合在实验量中的投影。
\end{derivation}

\begin{derivation}[从\eqnrefs{eq:heat-kernel-a4-flat-dp11}到\eqnrefs{eq:qcd-glue-ghost-uv-coefficient-dp10}]
纯 Yang--Mills 一圈系数中的 $11/3$ 来自规范矢量涨落与 Faddeev--Popov 鬼场对同一个背景场强平方项的不同热核贡献。其定量来源由规范固定后的二次算符分解给出。把完整规范联络分成经典背景场与量子涨落，

\begin{equation*}
 \mathscr C_\mu=\bar{\mathscr C}_\mu+Q_\mu.
\end{equation*}
背景 Feynman 规范取
\begin{equation*}
 \mathscr C_{\rm bg}^a[Q]=(\bar D^\mu Q_\mu)^a,
 \qquad\xi=1.
\end{equation*}
把 Yang--Mills 作用量与规范固定项展开到 $Q^2$；背景场满足经典运动方程时线性项消失，二次项可写为
\begin{equation*}
 S_E^{(2)}
 =\frac\hbar2\int\dd^4 x\,
 Q_\mu^a(\Delta_1)_{\mu\nu}^{ab}Q_\nu^b
 +\hbar\int\dd^4 x\,
 \bar c^a(\Delta_0)^{ab}c^b,
\end{equation*}
其中欧几里得算符为
\begin{align*}
 (\Delta_1)_{\mu\nu}
 &=-\bar D^2\delta_{\mu\nu}
 -2\ii g_s\bar{\mathcal F}_{\mu\nu}^aT_{\rm adj}^a,\\
 \Delta_0&=-\bar D^2.
\end{align*}
这里 $(T_{\rm adj}^a)_{bc}=-\ii f^{abc}$，所以
\begin{equation*}
 \Tr_{\rm adj}(T_{\rm adj}^aT_{\rm adj}^b)
 =C_A\delta^{ab}.
\end{equation*}
高斯积分给出纯规范扇区的一圈有效作用量
\begin{equation*}
 \Gamma_{g+\rm gh}^{(1)}
 =\frac12\Tr\ln\Delta_1-\Tr\ln\Delta_0.
\end{equation*}
第二项没有 $1/2$，因为 $c,\bar c$ 是彼此独立的复 Grassmann 场。

把这些算符逐项代入 \eqnrefs{eq:heat-kernel-a4-flat-dp11}。对背景伴随表示协变导数，丛曲率为
\begin{equation*}
 \Omega_{\rho\sigma}
 =[\bar D_\rho,\bar D_\sigma]
 =\ii g_s\bar{\mathcal F}_{\rho\sigma}^aT_{\rm adj}^a.
\end{equation*}
先做颜色迹：
\begin{align*}
 \Tr_{\rm adj}(\Omega_{\rho\sigma}\Omega_{\rho\sigma})
 &=(\ii g_s)^2
 \bar{\mathcal F}_{\rho\sigma}^a\bar{\mathcal F}_{\rho\sigma}^b
 \Tr_{\rm adj}(T^aT^b),\\
 &=-g_s^2C_A\bar{\mathcal F}_{\rho\sigma}^a
 \bar{\mathcal F}_{\rho\sigma}^a.
\end{align*}
矢量算符还需对四个 Lorentz 分量取迹，因此
\begin{equation*}
 \tr_{V\otimes\rm adj}\Omega^2
 =-4g_s^2C_A\bar{\mathcal F}^2.
\end{equation*}

与 Laplace 型标准形式 $\Delta=-(\nabla^2+E)$ 比较得到
\begin{equation*}
 (E_1)_{\mu\nu}
 =2\ii g_s\bar{\mathcal F}_{\mu\nu}^aT_{\rm adj}^a.
\end{equation*}
计算 $E_1^2$ 时，Lorentz 反对称性给
\begin{equation*}
 \bar{\mathcal F}_{\mu\nu}\bar{\mathcal F}_{\nu\mu}
 =-\bar{\mathcal F}_{\mu\nu}\bar{\mathcal F}_{\mu\nu}.
\end{equation*}
所以两个负号相消：
\begin{align*}
 \tr E_1^2
 &=(2\ii g_s)^2
 \bar{\mathcal F}_{\mu\nu}^a\bar{\mathcal F}_{\nu\mu}^b
 \Tr_{\rm adj}(T^aT^b),\\
 &=4g_s^2C_A\bar{\mathcal F}^2.
\end{align*}
因此矢量热核的对数局域系数为
\begin{align*}
 a_4(\Delta_1)
 &\supset
 \frac12\tr E_1^2+\frac1{12}\tr\Omega^2,\\
 &=\left(2-\frac13\right)g_s^2C_A\bar{\mathcal F}^2,\\
 &=\frac53g_s^2C_A\bar{\mathcal F}^2.
\end{align*}
对鬼场，$E_0=0$ 且没有矢量分量的重数：
\begin{equation*}
 a_4(\Delta_0)
 \supset-\frac1{12}g_s^2C_A\bar{\mathcal F}^2.
\end{equation*}
再乘上行列式对应的权重：
\begin{align*}
 \frac12a_4(\Delta_1)-a_4(\Delta_0)
 &=\frac12\frac53g_s^2C_A\bar{\mathcal F}^2
 -\left(-\frac1{12}g_s^2C_A\bar{\mathcal F}^2\right),\\
 &=\frac{11}{12}g_s^2C_A\bar{\mathcal F}^2.
\end{align*}
经典动能算符对应的拉格朗日量项为 $\bar{\mathcal F}^2/4$，因此把极点系数改写成相对于 $\frac14\bar{\mathcal F}^2$ 的系数需要乘以 $4$：
\begin{equation*}
 {4\times\frac{11}{12}C_A=\frac{11}{3}C_A.}
\end{equation*}
这就是 \eqnrefs{eq:qcd-glue-ghost-uv-coefficient-dp10} 的 $11C_A/3$。它不是某一张胶子圈图的独立数值，而是矢量涨落与 Faddeev--Popov 行列式在同一背景规范不变算符上的总和。
\end{derivation}

\begin{derivation}[从\eqnrefs{eq:reduced-loop-measure-si-r10}到\eqnrefs{eq:qcd-quark-loop-uv-coefficient-dp10}]
取质量为 $m_f$、属于颜色表示 $R_f$ 的一个 Dirac 费米子。颜色生成元满足
\begin{equation*}
 \Tr_{R_f}(T^aT^b)=T(R_f)\delta^{ab}.
\end{equation*}
采用统一的约化物理四动量圈积分记号，闭合费米子环的二点函数写成
\begin{equation*}
 \ii\Pi_{f,ab}^{\mu\nu}(q)
 =-g_s^2T(R_f)\delta_{ab}
 \int_\ell^{(d)}
 \frac{N^{\mu\nu}(\ell,q)}
 {D_0D_1},
 \qquad d=4-2\epsilon,
\end{equation*}
其中
\begin{align*}
 D_0&=\ell^2-m_f^2c^2+\ii0,\\
 D_1&=(\ell+q)^2-m_f^2c^2+\ii0,\\
 N^{\mu\nu}(\ell,q)
 &=\Tr\!\left[
 \gamma^\mu(\slashed\ell+m_fc)
 \gamma^\nu(\slashed\ell+\slashed q+m_fc)
 \right].
\end{align*}

先证明横向性，而不是在积分后凭 Lorentz 协变性猜结构。定义去掉整体 $\ii$ 的自由 Dirac 传播核
\begin{equation*}
 \mathsf S(\ell)
 \equiv\frac{\slashed\ell+m_fc}{D_0}
 =(\slashed\ell-m_fc)^{-1},
\end{equation*}
以及
\begin{equation*}
 \mathsf S(\ell+q)
 =\frac{\slashed\ell+\slashed q+m_fc}{D_1}.
\end{equation*}
于是
\begin{equation*}
 \slashed q
 =\mathsf S^{-1}(\ell+q)-\mathsf S^{-1}(\ell).
\end{equation*}
把 $q_\mu\gamma^\mu=\slashed q$ 代入循环迹，并只在迹内部使用循环性：
\begin{align*}
 &\Tr\!\left[
 \slashed q\,\mathsf S(\ell)\gamma^\nu\mathsf S(\ell+q)
 \right]\\
 &\quad=
 \Tr\!\left[
 \mathsf S^{-1}(\ell+q)\mathsf S(\ell)
 \gamma^\nu\mathsf S(\ell+q)
 \right]
 -\Tr\!\left[
 \mathsf S^{-1}(\ell)\mathsf S(\ell)
 \gamma^\nu\mathsf S(\ell+q)
 \right]\\
 &\quad=
 \Tr\!\left[
 \mathsf S(\ell)\gamma^\nu
 \mathsf S(\ell+q)\mathsf S^{-1}(\ell+q)
 \right]
 -\Tr\!\left[
 \gamma^\nu\mathsf S(\ell+q)
 \right]\\
 &\quad=
 \Tr\!\left[\mathsf S(\ell)\gamma^\nu\right]
 -\Tr\!\left[\gamma^\nu\mathsf S(\ell+q)\right].
\end{align*}
因此由起始定义逐项收缩可得
\begin{align*}
 \ii q_\mu\Pi_{f,ab}^{\mu\nu}(q)
 &=-g_s^2T(R_f)\delta_{ab}
 \int_\ell^{(d)}
 \Tr\!\left[
 \slashed q\,\mathsf S(\ell)\gamma^\nu\mathsf S(\ell+q)
 \right]\\
 &=-g_s^2T(R_f)\delta_{ab}
 \int_\ell^{(d)}
 \left\{
 \Tr[\mathsf S(\ell)\gamma^\nu]
 -\Tr[\gamma^\nu\mathsf S(\ell+q)]
 \right\}.
\end{align*}
第二项作 $r=\ell+q$ 平移。维数正规化保持积分测度平移不变，且 $\Tr[\mathsf S(r)\gamma^\nu]=\Tr[\gamma^\nu\mathsf S(r)]$，所以两项积分完全相消：
\begin{equation*}
 q_\mu\Pi_{f,ab}^{\mu\nu}(q)=0.
\end{equation*}
颜色对角性与 Lorentz 协变性先给出 $\Pi_{f,ab}^{\mu\nu}=\delta_{ab}[A_f(q^2)\eta^{\mu\nu}+B_f(q^2)q^\mu q^\nu]$；再与上式联立得到 $A_f=-q^2B_f$，因此存在唯一标量函数 $\Pi_f(q^2)$ 使
\begin{equation*}
 \Pi_{f,ab}^{\mu\nu}(q)
 =\delta_{ab}(q^2\eta^{\mu\nu}-q^\mu q^\nu)\Pi_f(q^2).
\end{equation*}

接着真正计算系数。Dirac 迹展开为
\begin{align*}
 N^{\mu\nu}
 &=\Tr(\gamma^\mu\slashed\ell\gamma^\nu\slashed\ell)
 +\Tr(\gamma^\mu\slashed\ell\gamma^\nu\slashed q)
 +m_f^2c^2\Tr(\gamma^\mu\gamma^\nu)\\
 &=4\left[
 \ell^\mu(\ell+q)^\nu
 +\ell^\nu(\ell+q)^\mu
 -\eta^{\mu\nu}\ell\cdot(\ell+q)
 +m_f^2c^2\eta^{\mu\nu}
 \right].
\end{align*}
奇数个 $\gamma$ 矩阵的两项为零。用 Feynman 参数
\begin{equation*}
 \frac1{D_0D_1}
 =\int_0^1\dd x\,
 \frac1{[xD_0+(1-x)D_1]^2}
\end{equation*}
并定义
\begin{equation*}
 r\equiv\ell+(1-x)q,
 \qquad
 \Delta_f(x)\equiv m_f^2c^2-x(1-x)q^2.
\end{equation*}
则
\begin{equation*}
 xD_0+(1-x)D_1=r^2-\Delta_f(x)+\ii0.
\end{equation*}
把
\begin{equation*}
 \ell=r-(1-x)q,
 \qquad
 \ell+q=r+xq
\end{equation*}
代入分子，保留关于 $r$ 的偶次项：
\begin{align*}
 \frac14N^{\mu\nu}_{\rm even}
 ={}&2r^\mu r^\nu-\eta^{\mu\nu}r^2
 -2x(1-x)q^\mu q^\nu\\
 &+\eta^{\mu\nu}
 \left[m_f^2c^2+x(1-x)q^2\right].
\end{align*}
所有关于 $r$ 的奇次项在对称积分中为零。

关键的张量积分不直接用“对称性可得”跳过。维数正规化下全导数积分为零：
\begin{equation*}
 0=\int_r^{(d)}
 \frac{\partial}{\partial r_\mu}
 \left[
 \frac{r_\nu}{r^2-\Delta_f+\ii0}
 \right].
\end{equation*}
展开导数得到
\begin{equation*}
 \int_r^{(d)}
 \frac{2r^\mu r^\nu}
 {(r^2-\Delta_f+\ii0)^2}
 =\eta^{\mu\nu}
 \int_r^{(d)}
 \frac1{r^2-\Delta_f+\ii0}.
\end{equation*}
另一方面
\begin{align*}
 \frac{r^2}{(r^2-\Delta_f)^2}
 &=\frac1{r^2-\Delta_f}
 +\frac{\Delta_f}{(r^2-\Delta_f)^2}.
\end{align*}
因此 $r^\mu r^\nu$ 与 $r^2$ 两项组合成
\begin{align*}
 &\int_r^{(d)}
 \frac{2r^\mu r^\nu-\eta^{\mu\nu}r^2}
 {(r^2-\Delta_f+\ii0)^2}\\
 &\qquad=-\eta^{\mu\nu}\Delta_f
 \int_r^{(d)}
 \frac1{(r^2-\Delta_f+\ii0)^2}.
\end{align*}
再加上分子中的显式质量与 $q^2$ 项，利用
\begin{equation*}
 -\Delta_f+m_f^2c^2+x(1-x)q^2
 =2x(1-x)q^2,
\end{equation*}
整个偶分子积分化为唯一横向组合
\begin{align*}
 \int_r^{(d)}\frac{N^{\mu\nu}_{\rm even}}
 {(r^2-\Delta_f+\ii0)^2}
 =8x(1-x)
 (q^2\eta^{\mu\nu}-q^\mu q^\nu)
 \int_r^{(d)}\frac1{(r^2-\Delta_f+\ii0)^2}.
\end{align*}

剩下的标量积分也直接算出。定义
\begin{equation*}
 I_2(\Delta)
 \equiv\frac1{\ii}\mu_{\rm DR}^{2\epsilon}
 \int\frac{\dd^d r}{(2\pi)^d}
 \frac1{(r^2-\Delta+\ii0)^2}.
\end{equation*}
Wick 转动或使用标准 Gamma 函数积分给
\begin{align*}
 I_2(\Delta)
 &=\frac{\mu_{\rm DR}^{2\epsilon}}
 {(4\pi)^{d/2}}
 \Gamma\!\left(2-\frac d2\right)
 (\Delta-\ii0)^{d/2-2}\\
 &=\frac1{(4\pi)^2}
 \left[
 \frac1{\bar\epsilon}
 -\ln\frac{\Delta-\ii0}{\mu_{\rm DR}^2}
 +O(\epsilon)
 \right],
\end{align*}
其中
\begin{equation*}
 \frac1{\bar\epsilon}
 \equiv\frac1\epsilon-\gamma_E+\ln4\pi.
\end{equation*}
因此紫外极点与 $m_f$ 和 $q^2$ 无关；这些尺度只进入有限对数。代回并使用
\begin{equation*}
 \int_0^1\dd x\,x(1-x)=\frac16,
\end{equation*}
得到单个 Dirac 表示对正文所采用背景动能紫外系数的贡献
\begin{equation*}
 \Pi_{f,ab}^{\mu\nu}(q)\Big|_{\rm UV}
 =\delta_{ab}(q^2\eta^{\mu\nu}-q^\mu q^\nu)
 \frac{g_s^2}{(4\pi)^2}
 \left[-\frac43T(R_f)\right]
 \frac1{\bar\epsilon}.
\end{equation*}
对所有活跃 Dirac 表示求和即得
\begin{equation*}
 -\frac43\sum_fT(R_f),
\end{equation*}
也就是正文\eqnrefs{eq:qcd-quark-loop-uv-coefficient-dp10} 的一般物质屏蔽项；QCD 基本表示 $R_f=F$ 只在最后一步代入。
\end{derivation}

\begin{derivation}[从\eqnrefs{eq:qcd-lagrangian-si-dp10}到\eqnrefs{eq:qcd-quark-gluon-vertex-dp10}]
把协变导数
\[
D_\mu=\partial_\mu+\ii g_s\mathscr C_\mu^aT^a
\]
代入夸克动能项，先把自由项与相互作用项严格分开：
\begin{align*}
\bar\psi_f\ii\hbar c\gamma^\mu D_\mu\psi_f
&=\bar\psi_f\ii\hbar c\gamma^\mu\partial_\mu\psi_f
+\bar\psi_f\ii\hbar c\gamma^\mu
(\ii g_s\mathscr C_\mu^aT^a)\psi_f\\
&=\bar\psi_f\ii\hbar c\gamma^\mu\partial_\mu\psi_f
-\hbar c g_s\bar\psi_f\gamma^\mu T^a\psi_f\,\mathscr C_\mu^a.
\end{align*}
因此三场相互作用拉氏量就是
\[
\Lag_{q\bar q g}
=-\hbar c g_s\bar\psi_{f,i}\gamma^\mu(T^a)_{ij}\psi_{f,j}\mathscr C_\mu^a.
\]
Dyson 展开的一阶项为
\begin{align*}
S_{q\bar q g}^{(1)}
&=\frac{\ii}{\hbar c}\int\dd^4x\,\Lag_{q\bar q g}(x)\\
&=-\ii g_s\int\dd^4x\,
\bar\psi_{f,i}(x)\gamma^\mu(T^a)_{ij}\psi_{f,j}(x)\mathscr C_\mu^a(x).
\end{align*}
把三条场都作 Fourier 展开，局域乘积中的时空积分产生整体四动量守恒因子
\[
\int\dd^4x\,
\exp\!\left[-\frac{\ii}{\hbar}
(p_1+p_2+k)\cdot x\right]
=(2\pi\hbar)^4\delta^{(4)}(p_1+p_2+k),
\]
这里按图规则约定三条四动量全部流入顶角。剥去所有图共有的动量守恒 $\delta$ 函数以后，剩下的局域矩阵系数恰为
\[
\boxed{-\ii g_s\gamma^\mu(T^a)_{ij}}.
\]
因此一个 QCD 顶角的三部分来源可以逐项追溯：$-\ii g_s$ 来自 Dyson 指数与协变导数中的耦合，$\gamma^\mu$ 来自 Dirac 电流，$T^a$ 来自颜色联络在基本表示上的生成元。图规则只是把这三个作用量系数压缩成一个可重复使用的记号。
\end{derivation}

\begin{derivation}[从\eqnrefs{eq:qqprime-amplitude-dp10}到\eqnrefs{eq:qqprime-M2-massive-dp115}]
从\eqnrefs{eq:qqprime-amplitude-dp10} 出发，初态自旋平均和颜色平均可分开写成
\begin{equation*}
 \overline{|\mathcal M|^2}
 =\frac{g_s^4}{t^2}\,
 \mathcal S_{\mu\nu}^{(1)}\mathcal S_{(2)}^{\mu\nu}\,\mathcal C,
\end{equation*}
其中
\begin{align*}
 \mathcal S_{\mu\nu}^{(1)}
 &=\frac12\Tr[(\slashed p_3+m_1c)\gamma_\mu(\slashed p_1+m_1c)\gamma_\nu],\\
 \mathcal S_{(2)}^{\mu\nu}
 &=\frac12\Tr[(\slashed p_4+m_2c)\gamma^\mu(\slashed p_2+m_2c)\gamma^\nu].
\end{align*}
颜色部分与质量无关：
\begin{align*}
 \mathcal C
 &=\frac1{N_c^2}\sum_{i,j,k,l}\sum_{a,b}
 (T^a)_{ki}(T^a)_{lj}(T^b)_{ik}(T^b)_{jl}\\
 &=\frac1{N_c^2}\sum_{a,b}\Tr(T^aT^b)\Tr(T^aT^b)\\
 &=\frac{T_F^2(N_c^2-1)}{N_c^2}
 =\frac{N_c^2-1}{4N_c^2}.
\end{align*}
Dirac 迹恒等式给
\begin{align*}
 \mathcal S_{\mu\nu}^{(1)}
 &=2\left[p_{3\mu}p_{1\nu}+p_{3\nu}p_{1\mu}
 -\eta_{\mu\nu}(p_1\!\cdot p_3-m_1^2c^2)\right],\\
 \mathcal S_{(2)}^{\mu\nu}
 &=2\left[p_4^{\mu}p_2^{\nu}+p_4^{\nu}p_2^{\mu}
 -\eta^{\mu\nu}(p_2\!\cdot p_4-m_2^2c^2)\right].
\end{align*}
先在四矢量标量积层面完成张量收缩。把两个方括号分别记为 $A_{\mu\nu}$ 与 $B^{\mu\nu}$，则逐项展开 $A_{\mu\nu}B^{\mu\nu}$ 后，所有
$(p_1\!\cdot p_3)(p_2\!\cdot p_4)$ 的纯迹项相消，得到
\begin{align*}
\mathcal S_{\mu\nu}^{(1)}\mathcal S_{(2)}^{\mu\nu}
=8\Big[&
(p_3\!\cdot p_4)(p_1\!\cdot p_2)
+(p_3\!\cdot p_2)(p_1\!\cdot p_4)\\
&-m_2^2c^2(p_1\!\cdot p_3)
-m_1^2c^2(p_2\!\cdot p_4)
+2m_1^2m_2^2c^4\Big].
\end{align*}
现在才把这些内积换成 Mandelstam 变量。除
\begin{align*}
 2p_1\!\cdot p_2&=s-\Sigma_{12},
&2p_1\!\cdot p_4&=\Sigma_{12}-u,\\
 2p_1\!\cdot p_3&=2m_1^2c^2-t,
&2p_2\!\cdot p_4&=2m_2^2c^2-t,
\end{align*}
外，由四动量守恒还可写成
\begin{equation*}
2p_3\!\cdot p_4=s-\Sigma_{12},
\qquad
2p_3\!\cdot p_2=\Sigma_{12}-u.
\end{equation*}
前两项因而给
\begin{equation*}
\frac14\left[(s-\Sigma_{12})^2+(u-\Sigma_{12})^2\right],
\end{equation*}
而三个显式质量项组合为
\begin{align*}
&-m_2^2c^2\left(m_1^2c^2-\frac t2\right)
-m_1^2c^2\left(m_2^2c^2-\frac t2\right)
+2m_1^2m_2^2c^4\\
&\qquad=\frac{t\Sigma_{12}}2.
\end{align*}
乘回最外层因子 $8$，得到
\begin{equation*}
 \mathcal S_{\mu\nu}^{(1)}\mathcal S_{(2)}^{\mu\nu}
 =2\left[(s-\Sigma_{12})^2+(u-\Sigma_{12})^2+2t\Sigma_{12}\right].
\end{equation*}
与颜色因子相乘即得\eqnrefs{eq:qqprime-M2-massive-dp115}。令 $m_1,m_2\to0$ 后，$\Sigma_{12}\to0$ 且 $s+t+u\to0$，从而得到\eqnrefs{eq:qqprime-M2-final-dp10}。
\end{derivation}

\begin{derivation}[从\eqnrefs{eq:qqbar-to-QQbar-amplitude-dp68}到\eqnrefs{eq:qqbar-to-QQbar-M2-dp10}]
同时保留初态质量 $m_q$ 与末态质量 $m_Q$。颜色与 Dirac 空间彼此独立，因此
\begin{equation*}
 \overline{|\mathcal M|^2}
 =\frac{g_s^4}{s^2}\,\mathcal C\,\mathcal S.
\end{equation*}
颜色平均为
\begin{align*}
\mathcal C
&=\frac1{N_c^2}
\sum_{i,j,k,l}\sum_{a,b}
(T^a)_{ji}(T^a)_{kl}
(T^b)_{ij}(T^b)_{lk}\\
&=\frac1{N_c^2}
\sum_{a,b}\Tr(T^aT^b)\Tr(T^aT^b)\\
&=\frac{T_F^2(N_c^2-1)}{N_c^2}
=\frac{N_c^2-1}{4N_c^2},
\end{align*}
其中基本表示取 $T_F=1/2$。

自旋平均后的两个 Dirac 迹为
\begin{align*}
 L^{\mu\nu}
 &\equiv\Tr[(\slashed p_2-m_qc)\gamma^\mu
 (\slashed p_1+m_qc)\gamma^\nu]\\
 &=4\left[p_2^\mu p_1^\nu+p_2^\nu p_1^\mu
 -\eta^{\mu\nu}(p_1\!\cdot p_2+m_q^2c^2)\right],\\
 H_{\mu\nu}
 &\equiv\Tr[(\slashed p_3+m_Qc)\gamma_\mu
 (\slashed p_4-m_Qc)\gamma_\nu]\\
 &=4\left[p_{3\mu}p_{4\nu}+p_{3\nu}p_{4\mu}
 -\eta_{\mu\nu}(p_3\!\cdot p_4+m_Q^2c^2)\right].
\end{align*}
于是
\begin{equation*}
 \mathcal S=\frac14L^{\mu\nu}H_{\mu\nu}.
\end{equation*}
先不换成 $s,t,u$，直接把两个对称张量逐项收缩。记
$a\equiv p_1\!\cdot p_2$、$b\equiv p_3\!\cdot p_4$，得到
\begin{align*}
\mathcal S
=8\Big[&
(p_1\!\cdot p_3)(p_2\!\cdot p_4)
+(p_1\!\cdot p_4)(p_2\!\cdot p_3)\\
&+m_Q^2c^2 a+m_q^2c^2 b
+2m_q^2m_Q^2c^4\Big].
\end{align*}
质量壳条件给
\begin{equation*}
a=\frac{s-2m_q^2c^2}{2},
\qquad
b=\frac{s-2m_Q^2c^2}{2}.
\end{equation*}
所以后三项不是独立的三种质量修正，而是精确合并成
\begin{align*}
&m_Q^2c^2\frac{s-2m_q^2c^2}{2}
+m_q^2c^2\frac{s-2m_Q^2c^2}{2}
+2m_q^2m_Q^2c^4\\
&\qquad=\frac{s}{2}(m_q^2+m_Q^2)c^2.
\end{align*}
因此
\begin{equation*}
 \mathcal S
 =8\left[
 (p_1\!\cdot p_3)(p_2\!\cdot p_4)
 +(p_1\!\cdot p_4)(p_2\!\cdot p_3)
 +\frac{s}{2}(m_q^2+m_Q^2)c^2
 \right].
\end{equation*}
令
\begin{equation*}
 \Sigma_{qQ}\equiv(m_q^2+m_Q^2)c^2.
\end{equation*}
由有限质量 Mandelstam 恒等式
\begin{align*}
 2p_1\!\cdot p_3&=\Sigma_{qQ}-t,
 &2p_2\!\cdot p_4&=\Sigma_{qQ}-t,\\
 2p_1\!\cdot p_4&=\Sigma_{qQ}-u,
 &2p_2\!\cdot p_3&=\Sigma_{qQ}-u,
\end{align*}
可得
\begin{equation*}
 \mathcal S
 =2\left[(t-\Sigma_{qQ})^2
 +(u-\Sigma_{qQ})^2
 +2s\Sigma_{qQ}\right].
\end{equation*}
乘回颜色因子得到
\begin{equation*}
 \boxed{
 \overline{|\mathcal M|^2}
 =g_s^4\frac{N_c^2-1}{2N_c^2}
 \frac{(t-\Sigma_{qQ})^2+(u-\Sigma_{qQ})^2+2s\Sigma_{qQ}}{s^2},}
\end{equation*}
即正文\eqnrefs{eq:qqbar-to-QQbar-M2-dp10}。外腿质量壳同时给
\begin{equation*}
 s+t+u=2\Sigma_{qQ}.
\end{equation*}
最后沿 $m_q,m_Q\to0$ 的连续质量壳极限得到
\begin{equation*}
 \lim_{m_q,m_Q\to0}\overline{|\mathcal M|^2}
 =g_s^4\frac{N_c^2-1}{2N_c^2}\frac{t^2+u^2}{s^2}.
\end{equation*}
无质量 Born 核因此是有限质量树级结果的严格连续极限。
\end{derivation}

\begin{derivation}[从\eqnrefs{eq:qcd-qqbar-gg-amplitude-sum-dp68}到\eqnrefs{eq:qcd-qqbar-gg-ward-dp32}]
把两条出射胶子动量记为 $k_1,k_2$，颜色指标为 $a,b$。将三张图共同的 $g_s^2$ 因子提出，并定义约化张量
\begin{equation*}
 \widetilde{\mathcal M}_x^{\mu\nu}
 \equiv g_s^{-2}\mathcal M_x^{\mu\nu},
 \qquad x=t,u,s.
\end{equation*}
对两张夸克交换图，费米子线结构严格写成
\begin{align*}
 \widetilde{\mathcal M}_t^{\mu\nu}
 &=\bar v(p_2)\gamma^\nu T^b
 S(p_1-k_1)\gamma^\mu T^a u(p_1),\\
 \widetilde{\mathcal M}_u^{\mu\nu}
 &=\bar v(p_2)\gamma^\mu T^a
 S(p_1-k_2)\gamma^\nu T^b u(p_1),
\end{align*}
其中 $S(q)$ 表示 Dirac 算符 $S^{-1}(q)=\slashed q-mc$ 的 Feynman 逆核。代数运算先在内部费米子分母非零的开运动学区域进行，所得恒等式再按统一的 $+\ii0$ 处方延拓。

把第一条胶子偏振替换为 $k_1^\mu$。由
\begin{equation*}
 \slashed k_1
 =S^{-1}(p_1)-S^{-1}(p_1-k_1)
\end{equation*}
以及在壳方程 $S^{-1}(p_1)u(p_1)=0$，得到
\begin{align*}
 S(p_1-k_1)\slashed k_1u(p_1)
 &=S(p_1-k_1)
 [S^{-1}(p_1)-S^{-1}(p_1-k_1)]u(p_1)\\
 &=-u(p_1).
\end{align*}
因此 $k_{1\mu}\widetilde{\mathcal M}_t^{\mu\nu}$ 留下一个没有内部传播子的费米子流，颜色顺序为 $T^bT^a$。

对另一种排序，利用四动量守恒 $p_1+p_2=k_1+k_2$，把 $k_1$ 写成靠近反费米子端的两个逆传播子之差，并令逆传播子向左作用到 $\bar v(p_2)$。由 $\bar v(p_2)(\slashed p_2+mc)=0$，得到同一个旋量流，颜色顺序为 $T^aT^b$，符号与上一项相反。两张约化夸克图的纵向收缩之和因而含有
\begin{equation*}
 T^aT^b-T^bT^a=[T^a,T^b]=\ii f^{abc}T^c
\end{equation*}
乘同一个费米子流。阿贝尔理论的生成元是普通数，所以这一对易子严格为零；QCD 的非零对易子要求三胶子图参与同一 Ward 恒等式。

$s$ 道图含夸克颜色流
\begin{equation*}
 J_c^\rho=\bar v(p_2)\gamma^\rho T^c u(p_1)
\end{equation*}
以及由 Yang--Mills 三次作用量得到的三胶子张量 $V_{\rho}{}^{\mu\nu}$。按“全部动量流入”的统一约定收缩 $k_{1\mu}V^{\rho\mu\nu}$，再使用 $k_1^2=k_2^2=0$、$k_2\cdot\varepsilon_2=0$ 与 $q=k_1+k_2$，得到
\begin{equation*}
 k_{1\mu}\widetilde{\mathcal M}_s^{\mu\nu}
 =-k_{1\mu}
 \left(\widetilde{\mathcal M}_t^{\mu\nu}
 +\widetilde{\mathcal M}_u^{\mu\nu}\right).
\end{equation*}
恢复共同因子 $g_s^2$ 后，三图满足
\begin{equation*}
 k_{1\mu}
 \left(\mathcal M_t^{\mu\nu}
 +\mathcal M_u^{\mu\nu}
 +\mathcal M_s^{\mu\nu}\right)=0.
\end{equation*}
交换 $1\leftrightarrow2$ 同样得到第二条 Ward 恒等式。三图抵消依次使用 Dirac 逆传播子的望远镜恒等式、颜色生成元的对易关系和由同一 Yang--Mills 作用量固定的三胶子顶角。
\end{derivation}

\begin{derivation}[从\eqnrefs{eq:qcd-bare-coupling-msbar-dp68}到\eqnrefs{eq:qcd-alpha-lambda-one-loop-dp10}]
重整化群方程不是额外加入 QCD 的动力学定律，而是同一个裸理论在改变参数定义尺度 \(\mu\) 时必须保持不变的条件。关键点是：在维数正规化的 \(d=4-2\epsilon\) 中，耦合本身先带有工程维数；这个 \(-\epsilon g_s\) 项与 \(Z_g\) 的紫外极点相乘，才留下四维中的有限运行。

从正文
\begin{equation*}
g_{s,B}=\mu^\epsilon Z_g(g_s,\epsilon)g_s,
\qquad
Z_g
=1-\frac{\beta_0}{2}
\frac{g_s^2}{(4\pi)^2}
\frac1{\bar\epsilon}
+\order(g_s^4)
\end{equation*}
出发。裸参数 \(g_{s,B}\) 不依赖人为的重整化尺度，因此
\begin{equation*}
0=\mu\frac{\dd}{\dd\mu}\ln g_{s,B}
=\epsilon
+\frac{\beta_g^{(d)}}{g_s}
+\beta_g^{(d)}\frac{\partial}{\partial g_s}\ln Z_g,
\end{equation*}
其中
\begin{equation*}
\beta_g^{(d)}(g_s)
\equiv\mu\frac{\dd g_s}{\dd\mu}
=-\epsilon g_s+\beta_g^{(4)}(g_s).
\end{equation*}
这里 \(-\epsilon g_s\) 是 \(d\) 维中的工程尺度项，而 \(\beta_g^{(4)}\) 是取四维极限后真正的量子运行。

只保留决定一圈有限项所需的最低阶。把 $\overline{\mathrm{MS}}$ 极点写成
\begin{equation*}
 \frac1{\bar\epsilon}
 \equiv
 \frac1\epsilon-\gamma_E+\ln(4\pi),
\end{equation*}
于是
\begin{align*}
\ln Z_g
&=-\frac{\beta_0}{2}
\frac{g_s^2}{(4\pi)^2}
\left[\frac1\epsilon-\gamma_E+\ln(4\pi)\right]
+\order(g_s^4),\\
\frac{\partial}{\partial g_s}\ln Z_g
&=-\frac{\beta_0 g_s}{(4\pi)^2}
\left[\frac1\epsilon-\gamma_E+\ln(4\pi)\right]
+\order(g_s^3).
\end{align*}
把 $\beta_g^{(d)}=-\epsilon g_s+\beta_g^{(4)}$ 代回裸尺度不变性。显式工程维数先给
\begin{equation*}
 \epsilon+\frac{-\epsilon g_s}{g_s}=0.
\end{equation*}
在剩下的乘积中，$\beta_g^{(4)}=\order(g_s^3)$ 乘上 $\partial_{g_s}\ln Z_g=\order(g_s/\epsilon)$ 只从两圈开始影响一圈方程，因此这一圈只需保留 $-\epsilon g_s$：
\begin{align*}
0
={}&\frac{\beta_g^{(4)}}{g_s}
+(-\epsilon g_s)
\left\{-\frac{\beta_0 g_s}{(4\pi)^2}
\left[\frac1\epsilon-\gamma_E+\ln(4\pi)\right]\right\}
+\order(g_s^4)+\order(\epsilon g_s^2)\\
={}&\frac{\beta_g^{(4)}}{g_s}
+\frac{\beta_0g_s^2}{(4\pi)^2}
\left[1-\epsilon\bigl(\gamma_E-\ln4\pi\bigr)\right]
+\order(g_s^4)+\order(\epsilon g_s^2).
\end{align*}
现在再取四维极限 $\epsilon\to0$，有限常数乘以 $\epsilon$ 的部分消失，而极点残数保留下来：
\begin{equation*}
0=\frac{\beta_g^{(4)}}{g_s}
+\frac{\beta_0g_s^2}{(4\pi)^2}
+\order(g_s^4).
\end{equation*}
因此
\begin{equation*}
\boxed{
\beta_g^{(4)}(g_s)
=-\frac{\beta_0}{16\pi^2}g_s^3
+\order(g_s^5).}
\end{equation*}
极点项 $1/\bar\epsilon$ 不能按普通代数量直接“约去”；有限的一圈运行来自严格极限
\begin{equation*}
 \lim_{\epsilon\to0}
 \epsilon\left(\frac1\epsilon-\gamma_E+\ln4\pi\right)=1.
\end{equation*}
这正是维数正规化中工程尺度项与紫外极点结合后留下有限反常尺度依赖的数学机制。

再用 \(\alpha_s=g_s^2/(4\pi)\)，链式法则给
\begin{align*}
\mu\frac{\dd\alpha_s}{\dd\mu}
&=\frac{g_s}{2\pi}
\mu\frac{\dd g_s}{\dd\mu}\\
&=-\frac{\beta_0}{2\pi}\alpha_s^2
+\order(\alpha_s^3),
\end{align*}
即正文\eqnrefs{eq:qcd-beta-alpha-one-loop-dp10}。

在活跃味数固定、\(\beta_0\) 为常数的尺度区间内，一圈方程可分离变量：
\begin{equation*}
\frac{\dd\alpha_s}{\alpha_s^2}
=-\frac{\beta_0}{2\pi}\dd\ln\mu.
\end{equation*}
从参考尺度 \(\mu_{\rm R0}\) 积到 \(\mathcal Q\)，
\begin{equation*}
\frac1{\alpha_s(\mathcal Q)}
-\frac1{\alpha_s(\mu_{\rm R0})}
=\frac{\beta_0}{2\pi}
\ln\frac{\mathcal Q}{\mu_{\rm R0}}.
\end{equation*}
把积分常数改写成具有动量量纲的 \(\Lambda_{\rm QCD}\)，定义
\begin{equation*}
\frac1{\alpha_s(\mathcal Q)}
\equiv
\frac{\beta_0}{4\pi}
\ln\frac{\mathcal Q^2}{\Lambda_{\rm QCD}^2},
\end{equation*}
便得到
\begin{equation*}
\boxed{
\alpha_s(\mathcal Q)
=\frac{4\pi}{\beta_0
\ln[\mathcal Q^2/\Lambda_{\rm QCD}^2]}.}
\end{equation*}
当 \(\beta_0>0\) 时，\(\mathcal Q\) 增大使分母增大，所以 \(\alpha_s\) 下降；这就是“渐近自由”在一圈阶的定量含义。反过来，\(\mathcal Q\) 接近 \(\Lambda_{\rm QCD}\) 时分母不再大，\(\alpha_s\) 失去小参数，说明微扰展开在红外失控；这一步本身并不等价于禁闭证明。
\end{derivation}

\begin{derivation}[从\eqnrefs{eq:qcd-r-born-sigma-si-dp12}到\eqnrefs{eq:qcd-r-nlo-final-dp12}]
固定一个无质量夸克味并令 $K=p_1+p_2$、$K^2=s$。以玻恩截面 $\sigma_{0,q}$ 归一化后，$O(\alpha_s)$ 只含两体虚修正与三体实辐射。

先算虚修正。重整化矢量流的在壳顶角可以写成树级电流乘标量形状因子。Feynman 规范下的一圈核为
\begin{equation*}
 \delta\Gamma^\mu
 =-\ii g_s^2 C_F\mu_{\rm DR}^{2\epsilon}
 \int\frac{\dd^d\ell}{(2\pi)^d}
 \frac{\gamma_\alpha(\slashed p_1-\slashed\ell)\gamma^\mu
 (-\slashed p_2-\slashed\ell)\gamma^\alpha}
 {[\ell^2+\ii0][(p_1-\ell)^2+\ii0][(p_2+\ell)^2+\ii0]},
 \qquad d=4-2\epsilon .
\end{equation*}
三分母用 $x+y+z=1$ 参数化：
\begin{equation*}
 \frac1{ABC}
 =2\int_0^1\dd x\int_0^{1-x}\dd y\,
 [xA+yB+(1-x-y)C]^{-3}.
\end{equation*}
令 $z=1-x-y$、$L=\ell-yp_1+zp_2$；利用 $p_1^2=p_2^2=0$ 与 $2p_1\!\cdot p_2=s$，分母成为
\begin{equation*}
 [L^2-yzs+\ii0]^3.
\end{equation*}
平移后的分子只需保留偶次 $L$。每出现一对圈动量就用
\begin{equation*}
 \int\dd^dL\,L^\rho L^\sigma F(L^2)
 =\frac{\eta^{\rho\sigma}}d\int\dd^dL\,L^2F(L^2)
\end{equation*}
约化，而剩余的 Dirac 链反复使用 $\gamma_\alpha\gamma_\rho\gamma^\alpha=(2-d)\gamma_\rho$ 与 $\{\gamma^\mu,\gamma^\nu\}=2\eta^{\mu\nu}$ 化简。再在左右分别作用 $\slashed p_1u(p_1)=0$、$\bar v(p_2)\slashed p_2=0$，所有允许的宇称偶矢量结构都退化为 $\bar v(p_2)\gamma^\mu u(p_1)$。所需标量积分只有
\begin{align*}
 I_3(s)
 &\equiv\mu_{\rm DR}^{2\epsilon}
 \int\frac{\dd^dL}{(2\pi)^d}\frac1{[L^2-yzs+\ii0]^3}\\
 &=-\frac{\ii}{(4\pi)^{2-\epsilon}}
 \frac{\Gamma(1+\epsilon)}{2}
 (yzs-\ii0)^{-1-\epsilon},
\end{align*}
以及由 $L^2=[L^2-yzs]+yzs$ 化成的二分母积分。把 $x,y$ 的多项式逐项积分后，空间样形状因子为
\begin{equation*}
 F_{q,\rm SL}^{(1)}
 =C_F\frac{\alpha_s}{4\pi}\,H(\epsilon)
 \left[-\frac2{\epsilon^2}-\frac3\epsilon-8+\order(\epsilon)\right],
 \qquad H(0)=1,
\end{equation*}
其中 $H(\epsilon)$ 收集共同的 $\overline{\rm MS}$ 尺度与 Gamma 函数因子。物理湮灭区是时间样 $s>0$，解析延拓
\begin{equation*}
 (-s-\ii0)^{-\epsilon}=s^{-\epsilon}\ee^{\ii\pi\epsilon},
 \qquad
 \cos(\pi\epsilon)=1-\frac{\pi^2\epsilon^2}{2}+\order(\epsilon^4)
\end{equation*}
使玻恩振幅与一圈振幅的干涉变为
\begin{equation*}
 \frac{\sigma_V}{\sigma_{0,q}}
 =C_F\frac{\alpha_s}{2\pi}H(\epsilon)
 \left[-\frac2{\epsilon^2}-\frac3\epsilon-8+\pi^2+\order(\epsilon)\right].
\end{equation*}
这里的双极点来自软--共线重叠区，单极点来自纯共线区；紫外极点已经包含在重整化定义中而被减去。

再算实辐射。取 $x_i=2k_i\!\cdot K/s$，则 $x_1+x_2+x_3=2$。定义三个非负变量
\begin{equation*}
 u\equiv1-x_1=\frac{2k_2\!\cdot k_3}{s},\qquad
 v\equiv1-x_2=\frac{2k_1\!\cdot k_3}{s},\qquad
 w\equiv1-x_3=\frac{2k_1\!\cdot k_2}{s},
 \qquad u+v+w=1.
\end{equation*}
$d$ 维三体相空间的角积分给出 $(uvw)^{-\epsilon}$；两张实辐射图的自旋--颜色求和则给出
\begin{equation*}
 \frac{x_1^2+x_2^2-\epsilon x_3^2}{(1-x_1)(1-x_2)}
 =\frac{2w+(1-\epsilon)(u^2+v^2)-2\epsilon uv}{uv}.
\end{equation*}
因此除去与虚图相同的公共因子 $C_F\alpha_sH(\epsilon)/(2\pi)$ 后，真正需要计算的是二维单纯形积分
\begin{align*}
 I_R(\epsilon)
 &=\int_0^1\dd u\int_0^{1-u}\dd v\,
 u^{-1-\epsilon}v^{-1-\epsilon}w^{-\epsilon}
 \bigl[2w+(1-\epsilon)(u^2+v^2)-2\epsilon uv\bigr].
\end{align*}
每一项都由 Dirichlet 积分
\begin{equation*}
 \int_{u,v,w\ge0}\dd u\,\dd v\,
 \delta(1-u-v-w)u^{a-1}v^{b-1}w^{c-1}
 =\frac{\Gamma(a)\Gamma(b)\Gamma(c)}{\Gamma(a+b+c)}
\end{equation*}
直接计算。逐项代入得
\begin{align*}
 I_R(\epsilon)
 &=2D(0,0,1)+(1-\epsilon)D(2,0,0)
 +(1-\epsilon)D(0,2,0)-2\epsilon D(1,1,0),
\end{align*}
其中
\begin{equation*}
 D(r,s,t)
 \equiv
 \frac{\Gamma(r-\epsilon)\Gamma(s-\epsilon)\Gamma(t+1-\epsilon)}
 {\Gamma(r+s+t+1-3\epsilon)}.
\end{equation*}
利用 $\Gamma(1-z)=-z\Gamma(-z)$ 与 $\Gamma(z+1)=z\Gamma(z)$，六项可以合并成
\begin{equation*}
 I_R(\epsilon)
 =-2(2\epsilon-1)(\epsilon^2-2\epsilon+2)
 \frac{\Gamma(-\epsilon)^2\Gamma(1-\epsilon)}
 {\Gamma(3-3\epsilon)}.
\end{equation*}
在 $\epsilon=0$ 附近展开 Gamma 函数，得到
\begin{equation*}
 I_R(\epsilon)
 =\frac2{\epsilon^2}+\frac3\epsilon
 +\frac{19}{2}-\pi^2+\order(\epsilon).
\end{equation*}
所以
\begin{equation*}
 \frac{\sigma_R}{\sigma_{0,q}}
 =C_F\frac{\alpha_s}{2\pi}H(\epsilon)
 \left[\frac2{\epsilon^2}+\frac3\epsilon
 +\frac{19}{2}-\pi^2+\order(\epsilon)\right].
\end{equation*}

现在实、虚两项逐阶相加：
\begin{align*}
 \frac{\sigma_R+\sigma_V}{\sigma_{0,q}}
 &=C_F\frac{\alpha_s}{2\pi}H(\epsilon)
 \left[
 \left(\frac2{\epsilon^2}-\frac2{\epsilon^2}\right)
 +\left(\frac3\epsilon-\frac3\epsilon\right)
 +\frac{19}{2}-8\right]\\
 &=C_F\frac{\alpha_s}{2\pi}\frac32+\order(\epsilon,\alpha_s^2).
\end{align*}
取 $C_F=4/3$ 得
\begin{equation*}
 \sigma_q^{\rm NLO}
 =\sigma_{0,q}\left[1+\frac{\alpha_s}{\pi}+\order(\alpha_s^2)\right].
\end{equation*}
最后对所有活跃夸克味求和并除以 $\mu^+\mu^-$ 玻恩截面，即得到正文\eqnrefs{eq:qcd-r-nlo-final-dp12}。
\end{derivation}

\begin{derivation}[从\eqnrefs{eq:wilson-gauge-action-dp11}到\eqnrefs{eq:wilson-continuum-limit-dp11}]
把格元中心记为 $x_c=x+a(\hat\mu+\hat\nu)/2$。
将中心处场强沿格边平行移回 $x$ 后，闭合路径展开为

\begin{equation*}
 U_{\mu\nu}(x)
 =\exp\!\left[\ii g_sa^2\mathcal F_{\mu\nu}(x_c)+\order(a^4)\right].
\end{equation*}
展开指数并使用 $\Tr T^a=0$ 与 $\Tr(T^aT^b)=\delta^{ab}/2$，有
\begin{align*}
 \Re\Tr U_{\mu\nu}
 &=N_c-
 \frac{g_s^2a^4}{2}
 \Tr\!\left(\mathcal F_{\mu\nu}(x_c)\mathcal F_{\mu\nu}(x_c)\right)
 +\order(a^6),\\
 &=N_c-
 \frac{g_s^2a^4}{4}
 \mathcal F_{\mu\nu}^a(x_c)\mathcal F_{\mu\nu}^a(x_c)
 +\order(a^6).
\end{align*}
若改在角点 $x$ 表示，同一个结果为
\begin{align*}
\Re\Tr U_{\mu\nu}(x)
=N_c-\frac{g_s^2a^4}{4}\Big[
\mathcal F_{\mu\nu}^a(x)\mathcal F_{\mu\nu}^a(x)
+\frac a2(\partial_\mu+\partial_\nu)
(\mathcal F_{\mu\nu}^a\mathcal F_{\mu\nu}^a)(x)\Big]
+\order(a^6).
\end{align*}
因此单个角点展开不能把余项直接写成 $\order(a^6)$；
$a^5$ 项要在周期边界或边界贡献可忽略的整体格点和中消去。
采用中心 Riemann 和再代入\eqnrefs{eq:wilson-gauge-action-dp11}，
并用 $\beta_{\rm lat}=2N_c/g_s^2$，得到
\begin{equation*}
 \frac{S_W}{\hbar}
 =\sum_xa^4\sum_{\mu<\nu}
 \frac12\mathcal F_{\mu\nu}^a\mathcal F_{\mu\nu}^a
 +\order(a^2).
\end{equation*}
反对称性给
\begin{equation*}
 \sum_{\mu<\nu}\mathcal F_{\mu\nu}^2
 =\frac12\sum_{\mu,\nu}\mathcal F_{\mu\nu}^2,
\end{equation*}
而 Riemann 和在 $a\to0$ 时变成四维积分，因此得到相应的欧几里得 Yang--Mills 连续极限。
\end{derivation}



\begin{derivation}[从\eqnrefs{eq:qcd-dirac-spectral-density-dp68}到\eqnrefs{eq:banks-casher-dp11}]
正文采用欧几里得 Dirac 算符的反厄米约定
\begin{equation*}
 \slashed D_E\phi_n=\ii\lambda_n\phi_n,
 \qquad \lambda_n\in\mathbb R,
 \qquad
 \int_{V_4}\dd^4 x\,\phi_m^\dagger\phi_n=\delta_{mn}.
\end{equation*}
无质量算符与 $\gamma^5$ 反对易：
\begin{equation*}
 \{\gamma^5,\slashed D_E\}=0.
\end{equation*}
因此若 $\lambda_n\neq0$，则
\begin{equation*}
 \slashed D_E(\gamma^5\phi_n)
 =-\ii\lambda_n(\gamma^5\phi_n),
\end{equation*}
非零谱严格成 $\pm\lambda$ 对。加入小而正的夸克质量，欧几里得传播子为
\begin{equation*}
 S_E=(\slashed D_E+mc/\hbar)^{-1},
\end{equation*}
因此单味凝聚由同点传播子的迹给出
\begin{equation*}
 -\langle\bar\psi\psi\rangle_m
 =\frac1{V_4}
 \left\langle\Tr\frac1{\slashed D_E+mc/\hbar}\right\rangle.
\end{equation*}
在 $\phi_n$ 本征基中插入完备关系：
\begin{equation*}
 \Tr\frac1{\slashed D_E+mc/\hbar}
 =\sum_n\frac1{\ii\lambda_n+mc/\hbar}.
\end{equation*}
对每一对非零本征值 $\pm\lambda$，
\begin{align*}
 \frac1{\ii\lambda+mc/\hbar}
 +\frac1{-\ii\lambda+mc/\hbar}
 &=\frac{2mc/\hbar}{\lambda^2+(mc/\hbar)^2}.
\end{align*}
于是
\begin{equation*}
 -\langle\bar\psi\psi\rangle_m
 =\frac1{V_4}\left\langle
 \frac{N_0}{mc/\hbar}
 +\sum_{\lambda_n>0}
 \frac{2mc/\hbar}{\lambda_n^2+(mc/\hbar)^2}
 \right\rangle,
\end{equation*}
其中 $N_0$ 是严格零模数。若固定拓扑扇区且不存在体积广延的额外零模，$N_0=O(1)$；更一般地只需 $\langle N_0\rangle_m/V_4\to0$。所以先取 $V_4\to\infty$ 后零模的密度贡献消失；自发凝聚来自零点附近数目随体积成比例增长的近零模，而不是有限个拓扑零模。

按正文\eqnrefs{eq:qcd-dirac-spectral-density-dp68}，热力学极限把全谱离散求和变成
\begin{equation*}
 \lim_{V_4\to\infty}
 \frac1{V_4}\left\langle\sum_n f(\lambda_n)\right\rangle
 =\int_{-\infty}^{+\infty}\dd\lambda\,\rho_m(\lambda)f(\lambda).
\end{equation*}
由于 $\rho_m(-\lambda)=\rho_m(\lambda)$，传播子迹中的奇函数部分积分为零，只剩
\begin{equation*}
 -\langle\bar\psi\psi\rangle_m
 =\int_{-\infty}^{+\infty}\dd\lambda\,
 \rho_m(\lambda)
 \frac{mc/\hbar}{\lambda^2+(mc/\hbar)^2}.
\end{equation*}
令 $\epsilon=mc/\hbar$。先在同一紫外调节器下扣除凝聚的解析紫外项，随后令 $\rho_m$ 在零点附近局部一致趋于连续的 $\rho$。分布极限
\begin{equation*}
 \lim_{\epsilon\to0^+}
 \frac{\epsilon}{\lambda^2+\epsilon^2}
 =\pi\delta(\lambda)
\end{equation*}
在这些假设下给出
\begin{align*}
 -\lim_{m\to0^+}\lim_{V_4\to\infty}
 \langle\bar\psi\psi\rangle_m
 &=\pi\int_{-\infty}^{+\infty}\dd\lambda\,
 \rho(\lambda)\delta(\lambda)\\
 &=\pi\rho(0),
\end{align*}
即正文\eqnrefs{eq:banks-casher-dp11}。若有 $N_f$ 个简并轻味，则每个味都满足同一关系，而味求和的凝聚相应多出因子 $N_f$。若把极限顺序反过来，有限体积谱仍是离散的，除有限个拓扑零模外没有连续的零点谱密度；因此 $V_4\to\infty$ 必须先于 $m\to0^+$，这正是自发对称性破缺所需的热力学极限。
\end{derivation}

\begin{derivation}[从\eqnrefs{eq:qcd-chiral-lo-action-dp33}到\eqnrefs{eq:qcd-gmor-dp33}]
取两轻味同位旋极限，并把

\begin{equation*}
 U=\exp\!\left(\frac{\ii\pi}{F_\pi}\right),
 \qquad
 \pi=\pi^a\tau^a,
 \qquad
 \Tr(\tau^a\tau^b)=2\delta^{ab}.
\end{equation*}
展开到二次阶，
\begin{align*}
 U&=1+\frac{\ii\pi}{F_\pi}-\frac{\pi^2}{2F_\pi^2}+\order(\pi^3),\\
 U^\dagger&=1-\frac{\ii\pi}{F_\pi}-\frac{\pi^2}{2F_\pi^2}+\order(\pi^3).
\end{align*}
动能项因而变成
\begin{align*}
 \frac{F_\pi^2}{4(\hbar c)^2}
 \Tr(\partial_\mu U^\dagger\partial^\mu U)
 &=\frac1{4(\hbar c)^2}
 \Tr(\partial_\mu\pi\,\partial^\mu\pi)+\order(\pi^4),\\
 &=\frac1{2(\hbar c)^2}
 \partial_\mu\pi^a\partial^\mu\pi^a+\order(\pi^4),
\end{align*}
得到三个规范归一化的实标量低能模。

质量矩阵在同位旋极限中为 $M_q=\hat m c^2 I_2$，其中 $\hat m=(m_u+m_d)/2$。线性 \(\pi\) 介子项在 $U+U^\dagger$ 中相消，而二次项给
\begin{align*}
 \Tr[M_q(U^\dagger+U)]
 &=4\hat m c^2
 -\frac{\hat m c^2}{F_\pi^2}\Tr(\pi^2)+\order(\pi^4),\\
 &=4\hat m c^2
 -\frac{2\hat m c^2}{F_\pi^2}\pi^a\pi^a+\order(\pi^4).
\end{align*}
代回\eqnrefs{eq:qcd-chiral-lo-action-dp33}，二次质量项为
\begin{equation*}
 -\frac{B_0\hat m c^2}{(\hbar c)^4}\pi^a\pi^a.
\end{equation*}
与标准二次作用量中的
\begin{equation*}
 -\frac{m_\pi^2c^4}{2(\hbar c)^4}\pi^a\pi^a
\end{equation*}
比较，得到
\begin{equation*}
 m_\pi^2c^4=2B_0\hat m c^2=B_0(m_u+m_d)c^2.
\end{equation*}
最后使用已定义的低能常数匹配
\begin{equation*}
 B_0=-\frac{(\hbar c)^3\langle\bar q q\rangle}{F_\pi^2},
\end{equation*}
即可得到\eqnrefs{eq:qcd-gmor-dp33}。该关系的修正由更高阶轻夸克质量插入与手征圈图共同控制，其相对大小按\eqnrefs{eq:qcd-chiral-control-dp33} 的幂次计数组织。
\end{derivation}

\begin{derivation}[从\eqnrefs{eq:ee-qqg-current-dp68}到\eqnrefs{eq:ee-threejet-distribution-dp14}]
取 $K^2=s$，并定义
\begin{equation*}
s_{12}=2p_1\!\cdot p_2,\qquad
s_{1g}=2p_1\!\cdot k,\qquad
s_{2g}=2p_2\!\cdot k,\qquad
s=s_{12}+s_{1g}+s_{2g}.
\end{equation*}
无质量条件 $p_1^2=p_2^2=k^2=0$ 使两条内部夸克传播子分别为 $s_{1g}^{-1}$ 与 $s_{2g}^{-1}$。把正文的强子流写成
\begin{equation*}
J_g^{\mu\rho,a}
=g_s\,\bar u(p_1)T^a\mathcal A^{\mu\rho}v(p_2),
\end{equation*}
其中
\begin{equation*}
\mathcal A^{\mu\rho}
=\frac{\gamma^\rho(\slashed p_1+\slashed k)\gamma^\mu}{s_{1g}}
-\frac{\gamma^\mu(\slashed p_2+\slashed k)\gamma^\rho}{s_{2g}}.
\end{equation*}
先检验外胶子的 Ward 恒等式。把 $\varepsilon_\rho^*$ 换成 $k_\rho$，第一项用
\begin{equation*}
\slashed k(\slashed p_1+\slashed k)
=2k\!\cdot p_1-\slashed p_1\slashed k
\end{equation*}
并左乘 $\bar u(p_1)$，得到 $+\bar uT^a\gamma^\mu v$；第二项同理给 $-\bar uT^a\gamma^\mu v$，所以
\begin{equation*}
k_\rho J_g^{\mu\rho,a}=0.
\end{equation*}
因此胶子物理极化求和中的规范纵向项不贡献，可以用 $-g_{\rho\sigma}$；电磁流守恒同样允许对虚光子的横向部分使用 $-g_{\mu\nu}$。

把三体结果除以同一虚光子产生 $q\bar q$ 的 Born 张量，颜色部分只留下
\begin{equation*}
\frac1{N_c}\sum_{a,i,j}(T^a)_{ij}(T^a)_{ji}
=\frac1{N_c}\sum_a\Tr(T^aT^a)=C_F.
\end{equation*}
自旋部分可统一写成一个迹：
\begin{equation*}
\mathcal T_3
=-g_{\mu\nu}(-g_{\rho\sigma})
\Tr\!\left[
\slashed p_1\mathcal A^{\mu\rho}
\slashed p_2\,\overline{\mathcal A}^{\nu\sigma}
\right],
\qquad
\overline{\mathcal A}^{\nu\sigma}
\equiv\gamma^0(\mathcal A^{\nu\sigma})^\dagger\gamma^0.
\end{equation*}
展开后共有两个平方项和两个干涉项。每个偶数 $\gamma$ 迹都用递推恒等式
\begin{equation*}
\Tr(\slashed a_1\cdots\slashed a_{2n})
=\sum_{j=2}^{2n}(-1)^j\,2a_1\!\cdot a_j\,
\Tr(\slashed a_2\cdots\widehat{\slashed a_j}\cdots\slashed a_{2n})
\end{equation*}
降到 $\Tr(\slashed a\slashed b)=4a\!\cdot b$。再用
\begin{equation*}
2p_1\!\cdot p_2=s_{12},\qquad
2p_1\!\cdot k=s_{1g},\qquad
2p_2\!\cdot k=s_{2g}
\end{equation*}
消去全部标量积，平方项和干涉项合并为标准的夸克--反夸克 antenna：
\begin{equation*}
\frac{\overline{|\mathcal M_{q\bar qg}|^2}}
{\overline{|\mathcal M_{q\bar q}|^2}}
=2g_s^2C_F
\left[
\frac{s_{1g}}{s\,s_{2g}}
+\frac{s_{2g}}{s\,s_{1g}}
+\frac{2s_{12}}{s_{1g}s_{2g}}
\right].
\end{equation*}
这一行仍保留了三个物理来源：前两项分别在 $p_2\parallel k$ 与 $p_1\parallel k$ 时产生共线极点，最后一项是两条辐射路径的相干干涉。

在质心系定义 $x_i=2E_i/E_{\rm cm}$，$E_{\rm cm}=c\sqrt s$。由
\begin{equation*}
(K-p_1)^2=s_{2g}=s(1-x_1),\qquad
(K-p_2)^2=s_{1g}=s(1-x_2),
\end{equation*}
以及 $x_1+x_2+x_3=2$，有
\begin{equation*}
s_{12}=s(1-x_3)=s(x_1+x_2-1).
\end{equation*}
于是 antenna 括号逐项化为
\begin{align*}
&s\left[
\frac{s_{1g}}{s\,s_{2g}}
+\frac{s_{2g}}{s\,s_{1g}}
+\frac{2s_{12}}{s_{1g}s_{2g}}
\right]\\
&\qquad=
\frac{(1-x_2)^2+(1-x_1)^2+2(x_1+x_2-1)}{(1-x_1)(1-x_2)}
=\frac{x_1^2+x_2^2}{(1-x_1)(1-x_2)}.
\end{align*}

最后计算相空间，而不是把 Jacobian 当作已知常数。对三个无质量末态，在去掉整体空间取向后，Dalitz 变量给出
\begin{equation*}
\dd\Phi_3
=\frac{1}{128\pi^3s}\,\dd s_{1g}\,\dd s_{2g}
=\frac{s}{128\pi^3}\,\dd x_1\,\dd x_2,
\end{equation*}
物理区域由 $s_{ij}>0$ 给出
\begin{equation*}
0<x_1<1,\qquad0<x_2<1,\qquad x_1+x_2>1.
\end{equation*}
同一总不变量质量下的无质量二体相空间积分为 $\Phi_2=1/(8\pi)$，所以
\begin{equation*}
\frac{\dd\Phi_3}{\Phi_2}
=\frac{s}{16\pi^2}\,\dd x_1\,\dd x_2.
\end{equation*}
把它与上面的振幅平方比相乘，并用 $g_s^2=4\pi\alpha_s$，得到
\begin{equation*}
\frac1{\sigma_0}
\frac{\dd^2\sigma}{\dd x_1\dd x_2}
=\frac{\alpha_s}{2\pi}C_F
\frac{x_1^2+x_2^2}{(1-x_1)(1-x_2)},
\end{equation*}
即正文\eqnrefs{eq:ee-threejet-distribution-dp14}。归一化到 $\sigma_0$ 后，Born 电磁流、入射通量与共同的二体归一化因子代数相消，剩余因子正是 QCD 辐射核与三体相空间 Jacobian。
\end{derivation}

\begin{derivation}[从\eqnrefs{eq:qcd-collinear-virtuality-dp68}到\eqnrefs{eq:qcd-lo-splitting-kernels-rest-dp83}]
把共线极限中出现的横向四矢量写成 $k_\perp^\mu$，并定义正的欧几里得横向大小
\begin{equation*}
 \kappa_\perp^2\equiv-k_\perp^2>0.
\end{equation*}
取两条类光参考方向 $P^\mu,n^\mu$，满足 $P^2=n^2=0$、$P\cdot n\neq0$。两个无质量子动量统一参数化为
\begin{align*}
 p^\mu
 &=zP^\mu+k_\perp^\mu
 +\frac{\kappa_\perp^2}{2zP\cdot n}\,n^\mu,\\
 k^\mu
 &=(1-z)P^\mu-k_\perp^\mu
 +\frac{\kappa_\perp^2}{2(1-z)P\cdot n}\,n^\mu,
\end{align*}
其中 $P\cdot k_\perp=n\cdot k_\perp=0$。直接平方可得
\begin{align*}
 p^2&=0,\qquad k^2=0,\\
 2p\cdot k
 &=\frac{\kappa_\perp^2}{z(1-z)},\\
 s_{pk}\equiv(p+k)^2
 &=\frac{\kappa_\perp^2}{z(1-z)},
\end{align*}
这与正文\eqnrefs{eq:qcd-collinear-virtuality-dp68} 相同。近在壳母线的传播子奇异因子 $1/s_{pk}$ 可以与子分裂顶角分离，而硬过程只解析到类光母方向 $P$。

先算 $g(P)\to q(p)\bar q(k)$。母胶子的物理偏振投影取轴向规范形式
\begin{equation*}
 d_{\mu\nu}(P,n)
 =-\eta_{\mu\nu}
 +\frac{P_\mu n_\nu+n_\mu P_\nu}{P\cdot n},
 \qquad
 P^\mu d_{\mu\nu}=n^\mu d_{\mu\nu}=0.
\end{equation*}
四维中母胶子有两个物理偏振，因此未极化分裂要乘 $1/2$。颜色平均给
\begin{align*}
 \frac1{N_c^2-1}\sum_{a,i,j}
 (T^a)_{ij}(T^a)_{ji}
 &=\frac1{N_c^2-1}\sum_a\Tr(T^aT^a)\\
 &=T_F.
\end{align*}
末态自旋求和后，分裂顶角的 Lorentz 分子为
\begin{align*}
 \mathcal N_{q\bar q}
 &\equiv\frac{T_F}{2}
 d_{\mu\nu}(P,n)
 \Tr\!\left(\slashed p\gamma^\mu\slashed k\gamma^\nu\right)\\
 &=2T_F\,d_{\mu\nu}(P,n)
 \left[p^\mu k^\nu+p^\nu k^\mu
 -\eta^{\mu\nu}p\cdot k\right],
\end{align*}
其中使用
\begin{equation*}
 \Tr(\gamma^\alpha\gamma^\mu\gamma^\beta\gamma^\nu)
 =4\left(
 \eta^{\alpha\mu}\eta^{\beta\nu}
 -\eta^{\alpha\beta}\eta^{\mu\nu}
 +\eta^{\alpha\nu}\eta^{\mu\beta}
 \right).
\end{equation*}
先收缩 $-\eta_{\mu\nu}$ 部分：
\begin{align*}
 -\eta_{\mu\nu}
 \Tr(\slashed p\gamma^\mu\slashed k\gamma^\nu)
 &=8p\cdot k
 =\frac{4\kappa_\perp^2}{z(1-z)}.
\end{align*}
再收缩轴向规范部分。由 Sudakov 分解，
\begin{equation*}
 P\cdot p=\frac{\kappa_\perp^2}{2z},\qquad
 P\cdot k=\frac{\kappa_\perp^2}{2(1-z)},\qquad
 n\cdot p=zP\cdot n,\qquad
 n\cdot k=(1-z)P\cdot n.
\end{equation*}
因此
\begin{align*}
 &\frac{P_\mu n_\nu+n_\mu P_\nu}{P\cdot n}
 \Tr(\slashed p\gamma^\mu\slashed k\gamma^\nu)\\
 &\quad=\frac{8}{P\cdot n}
 \left[(P\cdot p)(n\cdot k)
 +(P\cdot k)(n\cdot p)
 -(P\cdot n)(p\cdot k)\right]\\
 &\quad=4\kappa_\perp^2
 \left[\frac{1-z}{z}+\frac{z}{1-z}
 -\frac1{z(1-z)}\right]\\
 &\quad=-8\kappa_\perp^2.
\end{align*}
两部分相加并乘母偏振平均 $1/2$ 后，
\begin{align*}
 \mathcal N_{q\bar q}
 &=\frac{T_F}{2}
 \left[
 \frac{4\kappa_\perp^2}{z(1-z)}
 -8\kappa_\perp^2
 \right]\\
 &=2T_F\frac{\kappa_\perp^2}{z(1-z)}
 \left[z^2+(1-z)^2\right].
\end{align*}
近共线母传播子给出 $1/s_{pk}$ 奇点。把完整平方振幅比值严格分解为奇异部分与余项，
\begin{equation*}
 \frac{\overline{|\mathcal M_{g\to q\bar q}|^2}}
 {\overline{|\mathcal M_{g}|^2}}
 \equiv
 \mathcal S_{g\to q\bar q}(z,\kappa_\perp)
 +\mathcal R_{\rm coll}(z,\kappa_\perp),
\end{equation*}
并用近共线极限定义奇异部分
\begin{align*}
 \mathcal S_{g\to q\bar q}
 &\equiv
 g_s^2\frac{\mathcal N_{q\bar q}}{s_{pk}^2}\\
 &=2g_s^2T_F
 \frac{z(1-z)}{\kappa_\perp^2}
 \left[z^2+(1-z)^2\right],
\end{align*}
其中余项满足 $\lim_{\kappa_\perp\to0}\kappa_\perp^2\mathcal R_{\rm coll}=0$。因此分裂核由完整振幅中唯一的 $1/\kappa_\perp^2$ 领先奇点系数确定。
相空间因子也从 Lorentz 不变递推式得到。把两个近共线子粒子先合并成虚母动量 $Q=p+k$，则
\begin{equation*}
 \dd\Phi_{n+1}
 =\dd\Phi_n(\ldots,Q)\,
 \frac{\dd s_{pk}}{2\pi}\,
 \dd\Phi_2(Q;p,k).
\end{equation*}
在 $Q$ 的瞬时静止系中，对两个无质量子粒子有
\begin{equation*}
 \dd\Phi_2
 =\frac1{32\pi^2}\dd\Omega.
\end{equation*}
把方位角积分掉，并用 $z=(1+\cos\theta)/2$，即 $\dd\cos\theta=2\dd z$，得到
\begin{equation*}
 \int\dd\varphi\,\dd\Phi_2
 =\frac1{8\pi}\dd z.
\end{equation*}
因此每个近共线分裂相对于硬 $n$ 体相空间多出
\begin{align*}
 \dd\Phi_{\rm coll}
 &=\frac{\dd s_{pk}}{2\pi}\,
 \frac{\dd z}{8\pi}\\
 &=\frac1{16\pi^2}\dd s_{pk}\,\dd z.
\end{align*}
由
\begin{equation*}
 s_{pk}=\frac{\kappa_\perp^2}{z(1-z)}
\end{equation*}
并在固定 $z$ 的相空间坐标中取微分，
\begin{equation*}
 \dd s_{pk}
 =\frac{\dd\kappa_\perp^2}{z(1-z)},
\end{equation*}
所以
\begin{equation*}
 \dd\Phi_{\rm coll}
 =\frac{1}{16\pi^2}
 \frac{\dd z\,\dd\kappa_\perp^2}{z(1-z)}.
\end{equation*}
于是
\begin{align*}
 \dd P_{g\to q\bar q}
 &=\frac{\overline{|\mathcal M_{g\to q\bar q}|^2}}
 {\overline{|\mathcal M_{g}|^2}}
 \dd\Phi_{\rm coll}\\
 &=\frac{\alpha_s}{2\pi}
 T_F\left[z^2+(1-z)^2\right]
 \frac{\dd\kappa_\perp^2}{\kappa_\perp^2}\,\dd z,
\end{align*}
其中 $g_s^2=4\pi\alpha_s$。因此
\begin{equation*}
 P_{qg}(z)=T_F\left[z^2+(1-z)^2\right].
\end{equation*}

再算 $g(P)\to g(p)g(k)$。选取二维欧几里得横向基 $e_i^\mu$，满足
\begin{equation*}
 e_i\cdot e_j=-\delta_{ij},
 \qquad e_i\cdot P=e_i\cdot n=0,
 \qquad k_\perp^\mu=\kappa_i e_i^\mu.
\end{equation*}
母胶子的物理偏振可取 $\epsilon_0^\mu=e_i^\mu$。两个子胶子的横向基矢必须补上沿 $n^\mu$ 的分量才能同时满足轴向规范条件和各自的质量壳横向条件。逐一解
\begin{equation*}
 n\cdot\epsilon_1=p\cdot\epsilon_1=0,
 \qquad
 n\cdot\epsilon_2=k\cdot\epsilon_2=0
\end{equation*}
得到
\begin{align*}
 \epsilon_1^\mu(a)
 &=e_a^\mu+\frac{\kappa_a}{zP\cdot n}\,n^\mu,\\
 \epsilon_2^\mu(b)
 &=e_b^\mu-\frac{\kappa_b}{(1-z)P\cdot n}\,n^\mu.
\end{align*}
例如
\begin{align*}
 p\cdot\epsilon_1(a)
 &=k_\perp\cdot e_a
 +\frac{\kappa_a}{zP\cdot n}(zP\cdot n)\\
 &=-\kappa_a+\kappa_a=0,
\end{align*}
第二个子胶子同理。

三胶子顶角全部动量取流入。令 $Q=p+k$ 为近在壳母动量，并取
\begin{equation*}
 Q_0=-Q,
 \qquad Q_1=p,
 \qquad Q_2=k,
 \qquad Q_0+Q_1+Q_2=0.
\end{equation*}
去掉共同的 $-g_sf^{abc}$ 后，Lorentz 张量是
\begin{equation*}
 V_{\rho\mu\nu}
 =\eta_{\rho\mu}(Q_0-Q_1)_\nu
 +\eta_{\mu\nu}(Q_1-Q_2)_\rho
 +\eta_{\nu\rho}(Q_2-Q_0)_\mu.
\end{equation*}
定义收缩后的横向分子
\begin{equation*}
 \mathcal V_{iab}
 \equiv
 \epsilon_0^\rho(i)\epsilon_1^\mu(a)\epsilon_2^\nu(b)
 V_{\rho\mu\nu}.
\end{equation*}
三个张量项分别计算。共线领先阶可用 $Q=P+O(\kappa_\perp^2)$。第一项给
\begin{align*}
 (\epsilon_0\cdot\epsilon_1)
 (Q_0-Q_1)\cdot\epsilon_2
 &=-\delta_{ia}
 \left[-(1+z)P-k_\perp\right]\cdot
 \left[e_b-\frac{\kappa_b}{(1-z)P\cdot n}n\right]\\
 &=-\frac{2\delta_{ia}\kappa_b}{1-z}.
\end{align*}
第二项给
\begin{align*}
 (\epsilon_1\cdot\epsilon_2)
 (Q_1-Q_2)\cdot\epsilon_0
 &=(-\delta_{ab})
 \left[(2z-1)P+2k_\perp\right]\cdot e_i\\
 &=2\delta_{ab}\kappa_i.
\end{align*}
第三项给
\begin{align*}
 (\epsilon_2\cdot\epsilon_0)
 (Q_2-Q_0)\cdot\epsilon_1
 &=-\delta_{bi}
 \left[(2-z)P-k_\perp\right]\cdot
 \left[e_a+\frac{\kappa_a}{zP\cdot n}n\right]\\
 &=-\frac{2\delta_{ib}\kappa_a}{z}.
\end{align*}
三项相加，才得到
\begin{equation*}
 \mathcal V_{iab}
 =2\left[
 \delta_{ab}\kappa_i
 -\frac{\delta_{ib}\kappa_a}{z}
 -\frac{\delta_{ia}\kappa_b}{1-z}
 \right]
 +O\!\left(\frac{\kappa_\perp^2}{Q_h}\right).
\end{equation*}

颜色平均为
\begin{align*}
 \frac1{N_c^2-1}\sum_{a,b,c}f^{abc}f^{abc}
 &=\frac{C_A\delta^{aa}}{N_c^2-1}
 =C_A.
\end{align*}
定义
\begin{equation*}
 B_{iab}
 \equiv
 \delta_{ab}\kappa_i
 -\frac{\delta_{ib}\kappa_a}{z}
 -\frac{\delta_{ia}\kappa_b}{1-z},
\end{equation*}
则对两个子偏振求和并对两个母偏振平均需要计算
\begin{equation*}
 \frac12\sum_{i,a,b}\mathcal V_{iab}\mathcal V_{iab}
 =2\sum_{i,a,b}B_{iab}B_{iab}.
\end{equation*}
六类平方与交叉项分别为
\begin{align*}
 \sum_{i,a,b}(\delta_{ab}\kappa_i)^2
 &=2\kappa_\perp^2,\\
 \sum_{i,a,b}\left(\frac{\delta_{ib}\kappa_a}{z}\right)^2
 &=\frac{2\kappa_\perp^2}{z^2},\\
 \sum_{i,a,b}\left(\frac{\delta_{ia}\kappa_b}{1-z}\right)^2
 &=\frac{2\kappa_\perp^2}{(1-z)^2},\\
 -\frac{2}{z}
 \sum_{i,a,b}\delta_{ab}\delta_{ib}\kappa_\ii\kappa_a
 &=-\frac{2\kappa_\perp^2}{z},\\
 -\frac{2}{1-z}
 \sum_{i,a,b}\delta_{ab}\delta_{ia}\kappa_\ii\kappa_b
 &=-\frac{2\kappa_\perp^2}{1-z},\\
 \frac{2}{z(1-z)}
 \sum_{i,a,b}\delta_{ib}\delta_{ia}\kappa_a\kappa_b
 &=\frac{2\kappa_\perp^2}{z(1-z)}.
\end{align*}
因此
\begin{align*}
 \sum_{i,a,b}B_{iab}B_{iab}
 &=2\kappa_\perp^2
 \left[
 1+\frac1{z^2}+\frac1{(1-z)^2}
 -\frac1z-\frac1{1-z}+\frac1{z(1-z)}
 \right]\\
 &=2\kappa_\perp^2
 \frac{(1-z+z^2)^2}{z^2(1-z)^2},
\end{align*}
从而
\begin{equation*}
 \frac12\sum_{i,a,b}\mathcal V_{iab}\mathcal V_{iab}
 =4\kappa_\perp^2
 \frac{(1-z+z^2)^2}{z^2(1-z)^2}.
\end{equation*}
除以母传播子虚度平方 $s_{pk}^2=\kappa_\perp^4/[z^2(1-z)^2]$，再乘同一个共线相空间，得到
\begin{align*}
 \dd P_{g\to gg}
 &=\frac{\alpha_s}{2\pi}
 \left\{2C_A
 \frac{(1-z+z^2)^2}{z(1-z)}\right\}
 \frac{\dd\kappa_\perp^2}{\kappa_\perp^2}\,\dd z.
\end{align*}
最后使用恒等式
\begin{align*}
 \frac{(1-z+z^2)^2}{z(1-z)}
 &=\frac{z}{1-z}+\frac{1-z}{z}+z(1-z)
\end{align*}
得到实辐射核
\begin{equation*}
 P_{gg}^{\rm real}(z)
 =2C_A\left[
 \frac{z}{1-z}+\frac{1-z}{z}+z(1-z)
 \right].
\end{equation*}

最后，$q\to qg$ 的实辐射已经在正文\eqnrefs{eq:q-to-qg-splitting-probability-dp10} 中推得。若现在把被观测子粒子改成胶子，令其携带 $z_g=1-z_q$，则
\begin{align*}
 P_{gq}(z_g)
 &=P_{qq}^{\rm real}(1-z_g)\\
 &=C_F\frac{1+(1-z_g)^2}{z_g}.
\end{align*}
因此三个实辐射核全部从同一共线母式得到。$P_{gg}$ 的 $z\to1$ 端点必须与虚修正共同解释为加号分布，而 $\delta(1-z)$ 系数由夸克数与总动量守恒固定；合并这些端点项后得到正文\eqnrefs{eq:qcd-lo-splitting-kernels-rest-dp83}。
\end{derivation}



\begin{derivation}[从\eqnrefs{eq:qcd-gluon-propagator-dp10}到\eqnrefs{eq:qcd-three-gluon-vertex-dp10}]
Feynman 规范 $\xi=1$ 下，把纯规范二次项与规范固定项分部积分后写成
\begin{equation*}
 \mathcal L_2
 =\frac{\hbar c}{2}\,
 \mathscr C_\mu^a\,\eta^{\mu\nu}\partial^2\mathscr C_\nu^a .
\end{equation*}
因此作用在规范场上的二次微分算符是
\begin{equation*}
 (\mathsf K\mathscr C)_\mu^a
 =\delta^{ab}\eta_{\mu\nu}\partial^2\mathscr C^{b\nu}.
\end{equation*}
采用
\begin{equation*}
 \mathscr C_\mu^a(x)
 =\int\frac{\dd^4p}{(2\pi\hbar)^4}
 \widetilde{\mathscr C}_\mu^a(p)
 \ee^{-\ii p\cdot x/\hbar},
\end{equation*}
有 $\partial^2\mapsto-p^2/\hbar^2$，故
\begin{equation*}
 \mathsf K_{\mu\nu}^{ab}(p)
 =-\frac{p^2}{\hbar^2}\delta^{ab}\eta_{\mu\nu}.
\end{equation*}
Feynman 边界条件把 $p^2$ 替换为 $p^2+\ii0^+$。定义逆核 $\mathsf K_F^{-1}$ 使
\begin{equation*}
 \mathsf K_{\mu\rho}^{ac}(p)
 (\mathsf K_F^{-1})^{cb\rho}{}_{\nu}(p)
 =\delta^{ab}\delta_\mu{}^\nu,
\end{equation*}
于是逐指标求逆得到
\begin{equation*}
 (\mathsf K_F^{-1})_{\mu\nu}^{ab}(p)
 =-\hbar^2\frac{\delta^{ab}\eta_{\mu\nu}}{p^2+\ii0^+}.
\end{equation*}
在统一的约化振幅约定中，玻色 Feynman 线等于 $(\ii/\hbar^2)\mathsf K_F^{-1}$，从而
\begin{equation*}
 D_{\mu\nu}^{ab}(p)
 =\frac{\ii}{\hbar^2}(\mathsf K_F^{-1})_{\mu\nu}^{ab}(p)
 =\frac{-\ii\delta^{ab}\eta_{\mu\nu}}{p^2+\ii0^+},
\end{equation*}
即正文\eqnrefs{eq:qcd-gluon-propagator-dp10}。

三胶子规则来自 Yang--Mills 作用量的三次部分
\begin{equation*}
 \mathcal L_{3g}
 =\hbar c g_s f^{abc}
 (\partial_\mu\mathscr C_\nu^a)
 \mathscr C^{b\mu}\mathscr C^{c\nu}.
\end{equation*}
令三条流入场分别为 $(a,\alpha,p)$、$(b,\beta,q)$、$(c,\gamma,r)$，并约定
\begin{equation*}
 p+q+r=0.
\end{equation*}
三次场都是同一个玻色规范场，因此对作用量取三次泛函导数时，导数场可以对应三条外腿中的任意一条。把所有排列写到统一的颜色次序 $f^{abc}$ 后，三个 Lorentz 系数分别为
\begin{align*}
 \mathscr V_1^{\alpha\beta\gamma}
 &=\eta^{\alpha\beta}(p-q)^\gamma,\\
 \mathscr V_2^{\alpha\beta\gamma}
 &=\eta^{\beta\gamma}(q-r)^\alpha,\\
 \mathscr V_3^{\alpha\beta\gamma}
 &=\eta^{\gamma\alpha}(r-p)^\beta.
\end{align*}
这里每一个动量差都来自两种外腿排列的差：交换相邻两个胶子场时，$f^{abc}$ 变号，而 Fourier 导数把相应流入动量带入同一个 Lorentz 缩并。于是三阶作用量系数给出完整约化顶角
\begin{equation*}
 \boxed{
 V_{3g}^{abc,\alpha\beta\gamma}(p,q,r)
 =-g_sf^{abc}
 \left[
 \eta^{\alpha\beta}(p-q)^\gamma
 +\eta^{\beta\gamma}(q-r)^\alpha
 +\eta^{\gamma\alpha}(r-p)^\beta
 \right],}
\end{equation*}
即正文\eqnrefs{eq:qcd-three-gluon-vertex-dp10}。交换任意两条完整外腿，例如 $(a,\alpha,p)\leftrightarrow(b,\beta,q)$，$f^{abc}$ 与方括号同时变号，因此完整顶角保持玻色交换对称。
\end{derivation}

\begin{derivation}[从\eqnrefs{eq:qcd-qg-amplitude-dp94}到\eqnrefs{eq:qcd-qg-M2-dp94}]
先取轻夸克极限 $m_q\to0$。这一极限下三张图的相对系数仍由规范不变作用量固定。单张图本身具有规范依赖，因此偏振求和先保留两种横向物理偏振；只有完整振幅相加并满足 Ward 恒等式以后，才可以使用相应的协变化简。

先固定质心系，使每个四动量和物理偏振都能逐分量检查。
令 $Q\equiv\sqrt{s}/2$，并取
\begin{align*}
 p_1^\mu&=Q(1,0,0,1),&
 k_1^\mu&=Q(1,0,0,-1),\\
 p_2^\mu&=Q(1,\sin\theta,0,\cos\theta),&
 k_2^\mu&=Q(1,-\sin\theta,0,-\cos\theta).
\end{align*}
于是
\begin{equation*}
 s=4Q^2,
 \qquad
 t=-2Q^2(1-\cos\theta),
 \qquad
 u=-2Q^2(1+\cos\theta),
 \qquad s+t+u=0.
\end{equation*}
记 $c_\theta\equiv\cos\theta$。入射胶子沿 $-z$ 方向，可选两条实的横向偏振
\begin{equation*}
 \varepsilon_{1x}^\mu=(0,1,0,0),
 \qquad
 \varepsilon_{1y}^\mu=(0,0,1,0),
\end{equation*}
而出射胶子沿 $(-\sin\theta,0,-\cos\theta)$，可选
\begin{equation*}
 \varepsilon_{2\parallel}^\mu=(0,-c_\theta,0,\sin\theta),
 \qquad
 \varepsilon_{2y}^\mu=(0,0,1,0).
\end{equation*}
逐项有 $k_i\cdot\varepsilon_{i\lambda}=0$ 与 $\varepsilon_{i\lambda}^2=-1$。因此这四种偏振组合恰好穷尽两条外胶子的物理态，不含时间样或纵向态。

三张图的 Dirac 核分别定义为
\begin{align*}
 \mathscr A_s^{\mu\nu}
 &=\gamma^\nu\frac{\slashed p_1+\slashed k_1}{s}\gamma^\mu,\\
 \mathscr A_u^{\mu\nu}
 &=\gamma^\mu\frac{\slashed p_1-\slashed k_2}{u}\gamma^\nu,\\
 \mathscr A_t^{\mu\nu}
 &=\frac1t\Gamma^{\rho\mu\nu}(p_1-p_2,k_1,-k_2)\gamma_\rho.
\end{align*}
这里第一条三胶子腿的动量必须是 $p_1-p_2=k_2-k_1$，因为前文三胶子顶角约定三条动量全部流入顶角：
\begin{equation*}
 (p_1-p_2)+k_1-k_2=0.
\end{equation*}
这也是正文\eqnrefs{eq:qcd-qg-amplitude-dp94} 中动量方向的直接检查。

接着计算六个独立的 Dirac--偏振收缩。
对 $x,y\in\{s,u,t\}$ 定义
\begin{equation*}
 K_{xy}(\theta)
 \equiv
 \sum_{r_1,r_2}
 \sum_{\lambda_1,\lambda_2}
 \Bigl[
 \bar u_{r_2}(p_2)\mathscr A_x^{\mu\nu}u_{r_1}(p_1)
 \varepsilon_{1\lambda_1,\mu}
 \varepsilon_{2\lambda_2,\nu}
 \Bigr]
 \Bigl[
 \bar u_{r_2}(p_2)\mathscr A_y^{\alpha\beta}u_{r_1}(p_1)
 \varepsilon_{1\lambda_1,\alpha}
 \varepsilon_{2\lambda_2,\beta}
 \Bigr]^* .
\end{equation*}
利用无质量自旋和
\begin{equation*}
 \sum_r u_r(p)\bar u_r(p)=\slashed p
\end{equation*}
后，每个 $K_{xy}$ 都只是 $4\times4$ Dirac 矩阵的有限迹。含三胶子顶角的项还可以先把 Lorentz 三指标降成一个普通 Dirac 矩阵。令 $q=p_1-p_2=k_2-k_1$，则
\begin{align*}
 &\varepsilon_1^\mu\varepsilon_2^\nu
 \Gamma^{\rho}{}_{\mu\nu}(q,k_1,-k_2)\\
 &\qquad={}
 -2(k_1\!\cdot\varepsilon_2)\varepsilon_1^\rho
 +(\varepsilon_1\!\cdot\varepsilon_2)(k_1+k_2)^\rho
 -2(k_2\!\cdot\varepsilon_1)\varepsilon_2^\rho.
\end{align*}
这一步直接来自三胶子顶角的三项，并只用了
$k_1\cdot\varepsilon_1=k_2\cdot\varepsilon_2=0$ 与 $q=k_2-k_1$。对上面选定的四种物理偏振，它进一步化为
\begin{align*}
 \mathscr A_t^{(x,y)}
 &=-\frac{2Q\sin\theta}{t}\slashed\varepsilon_y,
 &
 \mathscr A_t^{(y,\parallel)}
 &=-\frac{2Q\sin\theta}{t}\slashed\varepsilon_y,\\
 \mathscr A_t^{(y,y)}
 &=-\frac{\slashed k_1+\slashed k_2}{t},
 &
 \mathscr A_t^{(x,\parallel)}
 &=\frac{-2Q\sin\theta\slashed\varepsilon_{1x}
 +c_\theta(\slashed k_1+\slashed k_2)
 -2Q\sin\theta\slashed\varepsilon_{2\parallel}}{t}.
\end{align*}
因此 $t$ 道也已被化成普通的有限伽马矩阵链，不再剩下未计算的三胶子张量。作为第一项示范，$s$ 道平方变成
\begin{align*}
 K_{ss}
 ={}&\sum_{\lambda_1,\lambda_2}
 \Tr\!\left[
 \slashed p_2
 \slashed\varepsilon_2
 \frac{\slashed p_1+\slashed k_1}{s}
 \slashed\varepsilon_1
 \slashed p_1
 \slashed\varepsilon_1
 \frac{\slashed p_1+\slashed k_1}{s}
 \slashed\varepsilon_2
 \right].
\end{align*}
把 $p_1+k_1=(2Q,0,0,0)$ 与上面的四条物理偏振逐一代入，并反复使用
\begin{equation*}
 \{\gamma^\mu,\gamma^\nu\}=2\eta^{\mu\nu}I_4,
 \qquad
 \Tr(\slashed a\slashed b\slashed c\slashed d)
 =4\bigl(a\!\cdot b\,c\!\cdot d-a\!\cdot c\,b\!\cdot d+a\!\cdot d\,b\!\cdot c\bigr),
\end{equation*}
四种偏振组合相加给
\begin{equation*}
 K_{ss}=4(1+c_\theta).
\end{equation*}

再显式算一个含三胶子顶角的干涉项 $K_{st}$。定义 Dirac 共轭矩阵
\begin{equation*}
 \overline{\mathscr A}_t
 \equiv\gamma^0\mathscr A_t^\dagger\gamma^0,
\end{equation*}
则固定一组物理偏振 $(\lambda_1,\lambda_2)$ 时的自旋和为
\begin{equation*}
 K_{st}^{(\lambda_1,\lambda_2)}
 =\Tr\!\left[
 \slashed p_2\,
 \mathscr A_s^{(\lambda_1,\lambda_2)}
 \slashed p_1\,
 \overline{\mathscr A}_t^{(\lambda_1,\lambda_2)}
 \right].
\end{equation*}
把上一段已经化简出的四个 $\mathscr A_t$ 与
$\mathscr A_s=\slashed\varepsilon_2(\slashed p_1+\slashed k_1)\slashed\varepsilon_1/s$
逐一代入，四组偏振分别给
\begin{align*}
 K_{st}^{(x,\parallel)}
 &=\frac{4(1+c_\theta)}{1-c_\theta},&
 K_{st}^{(x,y)}
 &=2(1+c_\theta),\\
 K_{st}^{(y,\parallel)}
 &=2(1+c_\theta),&
 K_{st}^{(y,y)}
 &=\frac{4(1+c_\theta)}{1-c_\theta}.
\end{align*}
这里前、后两项中的 $1/(1-c_\theta)$ 来自 $t=-2Q^2(1-c_\theta)$；分子中的 $Q^2$ 则由两条外动量的 Dirac 迹提供，所以能量尺度消掉，只留下无量纲角函数。四项相加便得到
\begin{align*}
 K_{st}
 &=2(1+c_\theta)
 +2(1+c_\theta)
 +\frac{8(1+c_\theta)}{1-c_\theta}\\
 &=\frac{4(1+c_\theta)(3-c_\theta)}{1-c_\theta}.
\end{align*}
夸克传播子图与三胶子交换图的干涉因此归结为普通 Dirac 迹。物理干涉项还要乘颜色 Gram 元素 $G_{st}<0$，所以 $K_{st}$ 的符号本身不能判断最终干涉是增强还是相消。

剩余四个独立收缩使用同一套迹恒等式得到
\begin{align*}
 K_{uu}&=\frac{8(1+c_\theta^2)}{1+c_\theta},&
 K_{tt}&=\frac{8(1+c_\theta)(c_\theta^2-2c_\theta+5)}{(1-c_\theta)^2},\\
 K_{su}&=4(1-c_\theta),&
 K_{st}&=\frac{4(1+c_\theta)(3-c_\theta)}{1-c_\theta},&
 K_{ut}&=\frac{8(c_\theta^2-c_\theta+2)}{1-c_\theta}.
\end{align*}
由于 $K_{yx}=K_{xy}^*$，这六项已经覆盖三张图平方后的全部九个平方与干涉项。单个 $K_{xy}$ 依赖所选物理偏振基；与颜色结构相干组合并对完整物理振幅求和后，最终结果不依赖这一基底选择。

Lorentz--Dirac 收缩确定后，再独立计算颜色 Gram 矩阵。
把三张图的颜色矩阵记为
\begin{equation*}
 \mathcal C_s=T^bT^a,
 \qquad
 \mathcal C_u=T^aT^b,
 \qquad
 \mathcal C_t=\ii f^{abc}T^c=[T^a,T^b].
\end{equation*}
注意 $\mathcal C_t^\dagger=-\mathcal C_t$，因此干涉项的符号不能靠“颜色因子平方”猜测。定义
\begin{equation*}
 G_{xy}\equiv\sum_{a,b}\Tr_F(\mathcal C_x\mathcal C_y^\dagger).
\end{equation*}
由
\begin{equation*}
 \sum_aT^aT^a=C_FI,
 \qquad
 T^aT^bT^a=\left(C_F-\frac{C_A}{2}\right)T^b,
 \qquad
 \Tr(T^aT^b)=T_F\delta^{ab},
\end{equation*}
逐项得到
\begin{align*}
 G_{ss}=G_{uu}&=T_Fd_AC_F,\\
 G_{su}=G_{us}&=T_Fd_A\left(C_F-\frac{C_A}{2}\right),\\
 G_{tt}&=T_Fd_AC_A,\\
 G_{st}=G_{ts}&=-\frac12T_Fd_AC_A,\\
 G_{ut}=G_{tu}&=+\frac12T_Fd_AC_A,
\end{align*}
其中 $d_A=N_c^2-1$、$T_F=1/2$。例如
\begin{align*}
 G_{st}
 &=\sum_{a,b}\Tr\!\left[T^bT^a(T^bT^a-T^aT^b)\right]\\
 &=T_Fd_A\left(C_F-\frac{C_A}{2}\right)-T_Fd_AC_F
 =-\frac12T_Fd_AC_A,
\end{align*}
负的 $G_{st}$ 与正的 $G_{ut}$ 直接显示非阿贝尔交换图与两种夸克排序具有不同的颜色干涉符号。

最后先在振幅层相干求和，再实施初态平均。
未平均的完整平方振幅为
\begin{align*}
 \sum_{\rm spin,pol,color}|\mathcal M|^2
 =g_s^4\Bigl[&
 G_{ss}K_{ss}+G_{uu}K_{uu}+G_{tt}K_{tt}
 +2G_{su}K_{su}+2G_{st}K_{st}+2G_{ut}K_{ut}
 \Bigr].
\end{align*}
把上面六个 $K_{xy}$ 与六个颜色系数代入，并用
\begin{equation*}
 C_F=\frac{N_c^2-1}{2N_c},
 \qquad C_A=N_c,
 \qquad d_A=N_c^2-1,
\end{equation*}
整理成
\begin{equation*}
 \sum_{\rm spin,pol,color}|\mathcal M|^2
 =g_s^4\frac{d_A}{N_c}
 \frac{(c_\theta^2+2c_\theta+5)
 \left[N_c^2(c_\theta^2+2c_\theta+5)-(1-c_\theta)^2\right]}
 {(1-c_\theta)^2(1+c_\theta)}.
\end{equation*}
初态有 $2N_c$ 个夸克自旋--颜色态和 $2d_A$ 个胶子物理偏振--颜色态，所以平均因子是 $1/(4N_cd_A)$。再使用
\begin{align*}
 s^2+u^2&=4Q^4(c_\theta^2+2c_\theta+5),\\
 t^2&=4Q^4(1-c_\theta)^2,\\
 su&=-8Q^4(1+c_\theta),
\end{align*}
就得到
\begin{equation*}
 \overline{|\mathcal M(qg\to qg)|^2}
 =g_s^4(s^2+u^2)
 \left[
 \frac1{t^2}-\frac{C_F}{C_A}\frac1{su}
 \right],
\end{equation*}
即正文\eqnrefs{eq:qcd-qg-M2-dp94}。这里 $1/t^2$ 与 $1/(su)$ 是六类运动学收缩和六类颜色 Gram 元素相干组合后的结果；它们不能分别归属于某一张 Feynman 图的“概率”。
\end{derivation}

\begin{derivation}[从\eqnrefs{eq:qcd-gg-color-amplitude-dp94}到\eqnrefs{eq:qcd-gg-M2-dp94}]
这一过程若直接用 $-\eta_{\mu\nu}$ 展开四条外腿的偏振求和，会产生很长的 Lorentz 指标收缩，而且容易把“偏振代数”和“颜色代数”混在一起。这里改用物理螺旋度基底：先把每个非零螺旋度振幅算成三个颜色张量的线性组合，再单独完成一般 $SU(N_c)$ 的颜色求和。这样每一步都只处理一种数学结构。

先固定质心系运动学与外胶子的物理偏振。
令 $Q\equiv\sqrt{s}/2$，取度规 $\eta_{\mu\nu}=\operatorname{diag}(1,-1,-1,-1)$，并写
\begin{align*}
 k_1^\mu&=Q(1,0,0,1),&
 k_2^\mu&=Q(1,0,0,-1),\\
 k_3^\mu&=Q(1,\sin\theta,0,\cos\theta),&
 k_4^\mu&=Q(1,-\sin\theta,0,-\cos\theta).
\end{align*}
于是
\begin{equation*}
 t=-\frac{s}{2}(1-\cos\theta),
 \qquad
 u=-\frac{s}{2}(1+\cos\theta),
 \qquad s+t+u=0.
\end{equation*}
四个外胶子都只取横向物理偏振。以右螺旋度为例，可选
\begin{align*}
 \varepsilon_R^\mu(k_1)&=\frac1{\sqrt2}(0,1,\ii,0),&
 \varepsilon_R^\mu(k_2)&=\frac1{\sqrt2}(0,-1,\ii,0),\\
 \varepsilon_R^{*\mu}(k_3)&=\frac1{\sqrt2}(0,-\cos\theta,\ii,\sin\theta),&
 \varepsilon_R^{*\mu}(k_4)&=\frac1{\sqrt2}(0,\cos\theta,\ii,-\sin\theta).
\end{align*}
左螺旋度由上式作复共轭得到。逐一检查可得
\begin{equation*}
 k_i\cdot\varepsilon_{R,L}(k_i)=0,
 \qquad
 \varepsilon_{R,L}^*(k_i)\cdot\varepsilon_{R,L}(k_i)=-1,
\end{equation*}
所以这里从一开始就没有引入纵向或时间样偏振。

为缩短颜色指标，定义
\begin{equation*}
 \mathcal C_s\equiv f^{abe}f^{cde},
 \qquad
 \mathcal C_t\equiv f^{ace}f^{bde},
 \qquad
 \mathcal C_u\equiv f^{ade}f^{bce}.
\end{equation*}
Jacobi 恒等式给出
\begin{equation*}
 \mathcal C_s-\mathcal C_t+\mathcal C_u=0.
\end{equation*}

随后逐图计算独立的螺旋度振幅。
先算 $g_Rg_R\to g_Rg_R$。为避免与整体系数的无关相位纠缠，下表列出费曼规则给出的 $\ii\mathcal M$。三行交换图分别含两个三胶子顶角和一个 $1/s$、$1/t$、$1/u$ 传播子，最后一行为四胶子接触顶角。把上面的四个偏振矢量逐分量代入三胶子顶角
\begin{equation*}
 \Gamma^{\mu\nu\rho}(p,q,r)
 =\eta^{\mu\nu}(p-q)^\rho
 +\eta^{\nu\rho}(q-r)^\mu
 +\eta^{\rho\mu}(r-p)^\nu,
 \qquad p+q+r=0,
\end{equation*}
即可得到
\begin{center}
\begin{tabular}{c|c}
图 & $\ii\mathcal M(g_Rg_R\to g_Rg_R)/(\ii g_s^2)$\\ \hline
$s$ 道 & $-\mathcal C_s\cos\theta$\\[2pt]
$t$ 道 & $\displaystyle
 \mathcal C_t\frac{19+7\cos\theta-11\cos^2\theta+\cos^3\theta}
 {4(1-\cos\theta)}$\\[8pt]
$u$ 道 & $\displaystyle
 \mathcal C_u\frac{19-7\cos\theta-11\cos^2\theta-\cos^3\theta}
 {4(1+\cos\theta)}$\\[8pt]
四胶子接触项 & $\displaystyle
 -\mathcal C_s\cos\theta
 -\frac{\mathcal C_t}{4}(3+2\cos\theta-\cos^2\theta)
 -\frac{\mathcal C_u}{4}(3-2\cos\theta-\cos^2\theta)$
\end{tabular}
\end{center}
这里每一行只是有限个四维分量的乘加。以 $s$ 道为例，传播子分母为 $(k_1+k_2)^2=s$；两个三胶子顶角分别与
$\varepsilon_R(k_1)\varepsilon_R(k_2)$ 和
$\varepsilon_R^*(k_3)\varepsilon_R^*(k_4)$ 收缩。利用
$k_i\cdot\varepsilon_i=0$ 后，分子中的 $Q^2=s/4$ 恰好抵消传播子的 $1/s$，留下表中的 $-\mathcal C_s\cos\theta$。$t$、$u$ 道保留的 $1/(1\mp\cos\theta)$ 正是后来前向、后向交换极点的来源；接触项没有传播子，因此只给角度多项式。

把四行相加，并用 $\mathcal C_s=\mathcal C_t-\mathcal C_u$，分子多项式发生大规模相消：
\begin{align*}
 \frac{\ii\mathcal M_{RR\to RR}}{\ii g_s^2}
 &=4\left(
 \frac{\mathcal C_t}{1-\cos\theta}
 +\frac{\mathcal C_u}{1+\cos\theta}
 \right)\\
 &=-2\left(
 \mathcal C_t\frac{s}{t}
 +\mathcal C_u\frac{s}{u}
 \right).
\end{align*}
这一步显示了四胶子接触项的作用：它本身没有交换极点，却与三个交换图的多项式部分相消，使完整振幅只剩满足 Ward 恒等式与交叉结构的组合。

其余非零螺旋度振幅由同一分量计算或交叉变换得到。取三个独立通道可写成
\begin{align*}
 \mathcal M_{RR\to RR}
 &=-2g_s^2s\left(\frac{\mathcal C_t}{t}+\frac{\mathcal C_u}{u}\right),\\
 \mathcal M_{LR\to LR}
 &=2g_s^2u\left(\frac{\mathcal C_t}{t}+\frac{\mathcal C_s}{s}\right),\\
 \mathcal M_{LR\to RL}
 &=-2g_s^2t\left(\frac{\mathcal C_s}{s}-\frac{\mathcal C_u}{u}\right).
\end{align*}
整体相位不影响平方振幅。宇称变换给出另外三个等模振幅；其余螺旋度组合在树级为零。因此螺旋度求和只需把这三个独立结果平方后乘以 $2$；因子 $2$ 来自宇称相关的另外三个等模通道。

螺旋度核确定后，计算一般 $SU(N_c)$ 的颜色 Gram 矩阵。
记伴随表示维数
\begin{equation*}
 d_A=N_c^2-1.
\end{equation*}
首先，三个颜色张量的平方范数相同。例如
\begin{align*}
 \sum_{a,b,c,d}\mathcal C_s^2
 &=\sum_{e,e'}
 \left(\sum_{a,b}f^{abe}f^{ab e'}\right)
 \left(\sum_{c,d}f^{cde}f^{cd e'}\right)\\
 &=C_A^2\sum_{e,e'}\delta^{ee'}\delta^{ee'}
 =C_A^2d_A.
\end{align*}
令伴随表示生成元为 $(F^a)_{bc}=-\ii f^{abc}$。由
\begin{equation*}
 \Tr(F^aF^b)=C_A\delta^{ab},
 \qquad
 \sum_aF^aF^bF^a=\frac{C_A}{2}F^b,
\end{equation*}
可得一个代表性的交叉颜色积
\begin{align*}
 \sum_{a,b,c,d}\mathcal C_s\mathcal C_t
 &=\sum_{e,e'}\Tr(F^eF^{e'}F^eF^{e'})\\
 &=\frac{C_A}{2}\sum_{e'}\Tr(F^{e'}F^{e'})
 =\frac12C_A^2d_A.
\end{align*}
再结合 Jacobi 关系 $\mathcal C_s-\mathcal C_t+\mathcal C_u=0$ 及外颜色指标的置换对称性，得到完整 Gram 矩阵
\begin{equation*}
 G_{xy}\equiv\sum_{a,b,c,d}\mathcal C_x\mathcal C_y
 =C_A^2d_A
 \begin{pmatrix}
 1&\tfrac12&-\tfrac12\\
 \tfrac12&1&\tfrac12\\
 -\tfrac12&\tfrac12&1
 \end{pmatrix}_{xy},
 \qquad x,y\in\{s,t,u\}.
\end{equation*}
负的 $G_{su}$ 并不表示负概率；它只记录两个颜色张量在颜色向量空间中的夹角，最终概率由完整振幅的 Gram 二次型给出。

最后完成颜色求和、螺旋度求和与初态平均。
利用上面的 Gram 矩阵，第一个独立螺旋度通道给出
\begin{align*}
 \sum_{\rm color}|\mathcal M_{RR\to RR}|^2
 &=4g_s^4C_A^2d_A\,s^2
 \left(
 \frac1{t^2}+\frac1{u^2}+\frac1{tu}
 \right).
\end{align*}
同理
\begin{align*}
 \sum_{\rm color}|\mathcal M_{LR\to LR}|^2
 &=4g_s^4C_A^2d_A\,u^2
 \left(
 \frac1{t^2}+\frac1{s^2}+\frac1{ts}
 \right),\\
 \sum_{\rm color}|\mathcal M_{LR\to RL}|^2
 &=4g_s^4C_A^2d_A\,t^2
 \left(
 \frac1{s^2}+\frac1{u^2}+\frac1{su}
 \right).
\end{align*}
宇称伙伴再给相同贡献，因此全螺旋度、全颜色求和为
\begin{equation*}
 \sum_{\rm hel,color}|\mathcal M|^2
 =8g_s^4C_A^2d_A\,\mathcal S(s,t,u),
\end{equation*}
其中
\begin{align*}
 \mathcal S(s,t,u)={}&
 s^2\left(\frac1{t^2}+\frac1{u^2}+\frac1{tu}\right)
 +t^2\left(\frac1{s^2}+\frac1{u^2}+\frac1{su}\right)\\
 &+u^2\left(\frac1{s^2}+\frac1{t^2}+\frac1{st}\right).
\end{align*}
利用 $s+t+u=0$，例如
\begin{equation*}
 -\frac{2su}{t^2}
 =\frac{s^2+u^2}{t^2}-1,
\end{equation*}
以及循环置换，可把它化成
\begin{equation*}
 \mathcal S(s,t,u)
 =2\left(
 3-\frac{su}{t^2}-\frac{st}{u^2}-\frac{tu}{s^2}
 \right).
\end{equation*}
每个初态胶子有 $d_A$ 个颜色和两个物理螺旋度，所以初态简并度为 $4d_A^2$。最终
\begin{align*}
 \overline{|\mathcal M|^2}
 &=\frac{1}{4d_A^2}
 \sum_{\rm hel,color}|\mathcal M|^2\\
 &=\frac{4g_s^4C_A^2}{d_A}
 \left(
 3-\frac{su}{t^2}-\frac{st}{u^2}-\frac{tu}{s^2}
 \right),
\end{align*}
再代入 $d_A=N_c^2-1$，即得到正文\eqnrefs{eq:qcd-gg-M2-dp94}。
\end{derivation}

\begin{derivation}[从\eqnrefs{eq:qcd-identical-qq-amplitude-dp95}到\eqnrefs{eq:qcd-identical-qq-m2-dp95}]
把两张图写为
\begin{align*}
\mathcal M_t&=\frac{g_s^2}{t}
[\bar u_3\gamma^\mu T^a u_1]
[\bar u_4\gamma_\mu T^a u_2],\\
\mathcal M_u&=\frac{g_s^2}{u}
[\bar u_4\gamma^\mu T^a u_1]
[\bar u_3\gamma_\mu T^a u_2].
\end{align*}
相同末态费米子交换给 $\mathcal M=\mathcal M_t-\mathcal M_u$，所以
\begin{equation*}
\overline{|\mathcal M|^2}
=\overline{|\mathcal M_t|^2}
+\overline{|\mathcal M_u|^2}
-2\Re\overline{\mathcal M_t\mathcal M_u^*}.
\end{equation*}
前两项直接由不同味结果得到
\begin{equation*}
\overline{|\mathcal M_t|^2}
=g_s^4\frac{N_c^2-1}{2N_c^2}\frac{s^2+u^2}{t^2},
\qquad
\overline{|\mathcal M_u|^2}
=g_s^4\frac{N_c^2-1}{2N_c^2}\frac{s^2+t^2}{u^2}.
\end{equation*}
真正的新信息在交叉项。颜色平均可写成一个单迹：
\begin{align*}
\mathcal C_{tu}
&=\frac1{N_c^2}\sum_{a,b}\Tr(T^aT^bT^aT^b).
\end{align*}
利用
\begin{equation*}
T^aT^bT^a=\left(C_F-\frac{C_A}{2}\right)T^b
=-\frac{1}{2N_c}T^b,
\end{equation*}
以及 $\sum_b\Tr(T^bT^b)=T_F(N_c^2-1)$，得到
\begin{equation*}
\mathcal C_{tu}
=-\frac{N_c^2-1}{4N_c^3}.
\end{equation*}
自旋交叉项不能拆成两个独立迹，而是形成一个八 $\gamma$ 链。先按完备关系把四条外腿的自旋求和，注意复共轭振幅会反转第二条费米子链的矩阵次序，于是
\begin{equation*}
\mathcal S_{tu}
=\frac14\Tr\!\left(
\slashed p_3\gamma^\mu\slashed p_1\gamma^\nu
\slashed p_4\gamma_\mu\slashed p_2\gamma_\nu
\right).
\end{equation*}
这里不能把迹“按视觉上相邻的四个矩阵”拆开。先把中间五个 $\gamma$ 矩阵作为整体收缩。四维 Clifford 代数给出
\begin{equation*}
\gamma^\mu\gamma^\alpha\gamma^\nu\gamma^\beta\gamma_\mu
=-2\gamma^\beta\gamma^\nu\gamma^\alpha,
\end{equation*}
因此
\begin{align*}
\mathcal S_{tu}
&=-\frac12\Tr\!\left(
\slashed p_3\slashed p_4\gamma^\nu
\slashed p_1\slashed p_2\gamma_\nu
\right).
\end{align*}
第二次收缩只剩四个普通四矢量。利用
\begin{equation*}
\gamma^\nu\slashed a\slashed b\gamma_\nu=4a\!\cdot b,
\qquad
\Tr(\slashed a\slashed b)=4a\!\cdot b,
\end{equation*}
逐步得到
\begin{align*}
\mathcal S_{tu}
&=-\frac12\,4(p_1\!\cdot p_2)\Tr(\slashed p_3\slashed p_4)\\
&=-8(p_1\!\cdot p_2)(p_3\!\cdot p_4).
\end{align*}
对无质量 $2\to2$ 运动学，$2p_1\!\cdot p_2=s$ 且 $2p_3\!\cdot p_4=s$，故
\begin{equation*}
\boxed{\mathcal S_{tu}=-2s^2.}
\end{equation*}
这一步说明干涉项只留下 $s^2$ 并不是由 $s+t+u=0$ 猜出的，而是八 $\gamma$ 链在 Clifford 代数下真正收缩后的结果。
于是
\begin{equation*}
-2\Re\overline{\mathcal M_t\mathcal M_u^*}
=-2\frac{g_s^4}{tu}\,\mathcal C_{tu}\mathcal S_{tu}
=-g_s^4\frac{N_c^2-1}{N_c^3}\frac{s^2}{tu},
\end{equation*}
与两个平方项相加即得到\eqnrefs{eq:qcd-identical-qq-m2-dp95}。注意颜色迹本身为负，而 Fermi 交换又给振幅相对负号；最终干涉符号必须在这两层结构全部保留后才能确定。
\end{derivation}



\begin{derivation}[从\eqnrefs{eq:qcd-field-strength-matrix-dp10}到\eqnrefs{eq:yang-mills-expanded-3-4-gluon-dp10}]
纯规范场部分的三次与四次相互作用可以从场强平方直接展开。定义

\begin{equation*}
\mathcal F_{\mu\nu}^{a(0)}\equiv
\partial_\mu\mathscr C_\nu^a-\partial_\nu\mathscr C_\mu^a,
\qquad
X_{\mu\nu}^a\equiv f^{abc}\mathscr C_\mu^b\mathscr C_\nu^c,
\end{equation*}
则 $\mathcal F_{\mu\nu}^a=\mathcal F_{\mu\nu}^{a(0)}-g_sX_{\mu\nu}^a$。平方为
\begin{equation*}
\mathcal F_{\mu\nu}^a\mathcal F^{a\mu\nu}
=\mathcal F_{\mu\nu}^{a(0)}\mathcal F^{a(0)\mu\nu}
-2g_s\mathcal F_{\mu\nu}^{a(0)}X^{a\mu\nu}
+g_s^2X_{\mu\nu}^aX^{a\mu\nu}.
\end{equation*}
乘上 $-\hbar c/4$ 后，依次得到自由项、三次项与四次项。三次项的系数还需要利用场强的反对称结构整理：
\begin{align*}
\mathcal F_{\mu\nu}^{a(0)}X^{a\mu\nu}
&=f^{abc}(\partial_\mu\mathscr C_\nu^a-\partial_\nu\mathscr C_\mu^a)
\mathscr C^{b\mu}\mathscr C^{c\nu}.
\end{align*}
第二项交换哑 Lorentz 指标 $\mu\leftrightarrow\nu$ 后变为
\begin{equation*}
-f^{abc}(\partial_\mu\mathscr C_\nu^a)
\mathscr C^{b\nu}\mathscr C^{c\mu}.
\end{equation*}
再交换哑颜色指标 $b\leftrightarrow c$，并利用 $f^{acb}=-f^{abc}$，它变成与第一项相同的
\begin{equation*}
f^{abc}(\partial_\mu\mathscr C_\nu^a)
\mathscr C^{b\mu}\mathscr C^{c\nu}.
\end{equation*}
因此两项相加产生因子 $2$，三次拉氏量可写成
\begin{equation*}
\mathcal L_{3g}
=\hbar c g_s f^{abc}(\partial_\mu\mathscr C_\nu^a)
\mathscr C^{b\mu}\mathscr C^{c\nu},
\end{equation*}
这就是\eqnrefs{eq:yang-mills-expanded-3-4-gluon-dp10} 的交叉项。四次项为
\begin{equation*}
\mathcal L_{4g}
=-\frac{\hbar c g_s^2}{4}
 f^{abc}f^{ade}
 \mathscr C_\mu^b\mathscr C_\nu^c
 \mathscr C^{d\mu}\mathscr C^{e\nu}.
\end{equation*}

\end{derivation}





\begin{derivation}[有限靶质量的纵向投影]
前文 $F_L=F_2-2x_BF_1$ 是领先靶质量约定。若保留靶质量 $M_H$，从\eqnrefs{eq:dis-transverse-tensor-basis-dp68} 得到
\begin{equation*}
 \widetilde P^2=M_H^2c^2+\frac{(P\cdot q)^2}{\mathcal Q^2},\qquad
 \varepsilon_L^\mu=\frac{\widetilde P^\mu}{\sqrt{\widetilde P^2}},\qquad
 \varepsilon_L^2=1,\quad q\cdot\varepsilon_L=0.
\end{equation*}
类空虚光子正交子空间中的这个类时向量不等同于实光子的横向极化。其强子张量投影为
\begin{align*}
 W_L&=\varepsilon_L^\mu W_{\mu\nu}\varepsilon_L^\nu
 =-F_1+\frac{\widetilde P^2}{P\cdot q}F_2,\\
 2x_BW_L&=\left(1+\frac{4x_B^2M_H^2c^2}{\mathcal Q^2}\right)F_2-2x_BF_1.
\end{align*}
所以保留有限靶质量时，物理纵向组合应记为
\begin{equation}
 F_L^{\rm phys}\equiv
 \left(1+\frac{4x_B^2M_H^2c^2}{\mathcal Q^2}\right)F_2-2x_BF_1.
 \label{eq:dis-longitudinal-target-mass}
\end{equation}
当 $4x_B^2M_H^2c^2/\mathcal Q^2\ll1$ 时才恢复前述直线分离的领先幂形式。这个运动学修正与结构函数内部的高扭曲动力学修正不同：只修改纵向定义，并不能把全部有限 $\mathcal Q$ 效应一并纳入。
\end{derivation}

\begin{derivation}[从\eqnrefs{eq:qcd-soft-leg-factor-dp68}到\eqnrefs{eq:qcd-color-charge-conservation-amplitude-dp10}]
先取一条出射无质量夸克外腿。把不含额外软胶子的其余硬过程记为颜色--旋量对象 $\mathcal H_i$，则把四动量 $k$、颜色 $a$、偏振 $\varepsilon_\lambda$ 的胶子接到该外腿后，相关因子为
\begin{equation*}
 \mathcal M_{i}^{a}(k,\lambda)
 =g_s\,\bar u(p_i)\slashed\varepsilon_\lambda^*(k)T_i^a
 \frac{\slashed p_i+\slashed k}{(p_i+k)^2+\ii0}\,\mathcal H_i .
\end{equation*}
外腿与软胶子都在壳，$p_i^2=k^2=0$，所以传播子分母为
\begin{equation*}
 (p_i+k)^2=2p_i\cdot k.
\end{equation*}
分子中先处理不含 $k$ 的部分。由 Clifford 代数
\begin{equation*}
 \slashed\varepsilon^*\slashed p_i
 =2p_i\cdot\varepsilon^*-\slashed p_i\slashed\varepsilon^*
\end{equation*}
以及外腿 Dirac 方程 $\bar u(p_i)\slashed p_i=0$，有
\begin{align*}
 \bar u(p_i)\slashed\varepsilon^*(\slashed p_i+\slashed k)
 &=2p_i\cdot\varepsilon^*\,\bar u(p_i)
   +\bar u(p_i)\slashed\varepsilon^*\slashed k,\\
 \frac{\bar u\slashed\varepsilon^*\slashed k}{2p_i\cdot k}
 &=\order(k^0),
 \qquad
 \frac{2p_i\cdot\varepsilon^*\,\bar u}{2p_i\cdot k}
 =\frac{p_i\cdot\varepsilon^*}{p_i\cdot k}\bar u.
\end{align*}
因此随软能量 $\omega$ 最强的 $\omega^{-1}$ 项为
\begin{equation*}
 \mathcal M_i^{a}(k,\lambda)
 =g_sT_i^a\frac{p_i\cdot\varepsilon_\lambda^*}{p_i\cdot k}\,
 \mathcal M_n+\order(\omega^0).
\end{equation*}
入射外腿通过 crossing 写成全出射 convention 时会多出相反的颜色荷方向。统一记
\begin{equation*}
 \sigma_i^{\rm io}=
 \begin{cases}
 +1,&\text{出射外腿},\\
 -1,&\text{入射外腿},
 \end{cases}
\end{equation*}
即可得到正文\eqnrefs{eq:qcd-soft-leg-factor-dp68}。软胶子在领先阶不能分辨硬子图内部的短距离结构，故把所有外腿贡献相加：
\begin{equation*}
 |\mathcal M_{n+1}^{a}(k,\lambda)\rangle
 =g_s\sum_i\sigma_i^{\rm io}\mathbf T_i^a
 \frac{p_i\cdot\varepsilon_\lambda^*}{p_i\cdot k}
 |\mathcal M_n\rangle+\order(\omega^0),
\end{equation*}
即正文\eqnrefs{eq:qcd-soft-gluon-theorem-dp10}。

颜色荷守恒不是附加假设，而由同一振幅的 Ward 恒等式给出。把外胶子偏振替换为纵向方向 $\varepsilon_\lambda^\mu\to k^\mu$，领先软项变成
\begin{align*}
 \left.|\mathcal M_{n+1}^{a}\rangle\right|_{\varepsilon\to k}
 &=g_s\sum_i\sigma_i^{\rm io}\mathbf T_i^a
 \frac{p_i\cdot k}{p_i\cdot k}|\mathcal M_n\rangle+\order(k)\\
 &=g_s\left(\sum_i\sigma_i^{\rm io}\mathbf T_i^a\right)|\mathcal M_n\rangle+\order(k).
\end{align*}
物理外胶子振幅满足 $k_\mu\mathcal M^{a\mu}=0$。取 $k\to0$ 后，$\order(k)$ 项消失，于是
\begin{equation*}
 \left(\sum_i\sigma_i^{\rm io}\mathbf T_i^a\right)|\mathcal M_n\rangle=0,
\end{equation*}
得到正文\eqnrefs{eq:qcd-color-charge-conservation-amplitude-dp10}。因此软定理中的纵向规范依赖与硬振幅的总颜色荷恰好由同一 Ward 恒等式消去。
\end{derivation}

\begin{derivation}[从\eqnrefs{eq:qcd-collinear-virtuality-dp68}到\eqnrefs{eq:q-to-qg-splitting-probability-dp10}]
令类光母方向 $P^\mu$ 与辅助类光方向 $n^\mu$ 满足
\begin{equation*}
 P^2=n^2=0,
 \qquad P\cdot n\neq0,
\end{equation*}
并取横向矢量 $k_\perp\cdot P=k_\perp\cdot n=0$。记其欧几里得平方为
\begin{equation*}
 K^2\equiv-k_\perp^2>0.
\end{equation*}
使两个子粒子严格在壳的 Sudakov 参数化为
\begin{align*}
 p^\mu
 &=zP^\mu+k_\perp^\mu
 +\frac{K^2}{z}\frac{n^\mu}{2P\cdot n},\\
 k^\mu
 &=(1-z)P^\mu-k_\perp^\mu
 +\frac{K^2}{1-z}\frac{n^\mu}{2P\cdot n}.
\end{align*}
直接平方可得
\begin{align*}
 p^2
 &=-K^2+2z(P\cdot n)\frac{K^2}{2z(P\cdot n)}=0,\\
 k^2
 &=-K^2+2(1-z)(P\cdot n)
 \frac{K^2}{2(1-z)(P\cdot n)}=0.
\end{align*}
离壳母动量 $\widetilde P=p+k$ 因而为
\begin{equation*}
 \widetilde P^\mu
 =P^\mu+\frac{K^2}{z(1-z)}\frac{n^\mu}{2P\cdot n},
\end{equation*}
所以
\begin{equation*}
 \widetilde P^2=\frac{K^2}{z(1-z)},
\end{equation*}
即正文\eqnrefs{eq:qcd-collinear-virtuality-dp68}。

把分裂顶角连到任意硬子图 $\mathcal H$。共线极限下，与硬子图相连的旋量密度矩阵只需保留母方向 $\slashed P$。对子夸克自旋和胶子的两个物理偏振求和时，先在轴向规范 $n\cdot A=0$ 中选择两个横向基 $e_a^\mu$，满足
\begin{equation*}
 e_a\cdot P=e_a\cdot n=0,
 \qquad e_a\cdot e_b=-\delta_{ab}.
\end{equation*}
写 $k_\perp^\mu=\kappa_a e_a^\mu$。满足 $k\cdot\varepsilon_a=n\cdot\varepsilon_a=0$ 的物理偏振可取
\begin{equation*}
 \varepsilon_a^\mu(k)
 =e_a^\mu-\frac{\kappa_a}{(1-z)P\cdot n}n^\mu.
\end{equation*}
分裂部分的自旋密度矩阵为
\begin{equation*}
 \mathcal S
 =\sum_{a=1}^{2}
 \frac{\slashed{\widetilde P}\slashed\varepsilon_a^*\slashed p
 \slashed\varepsilon_a\slashed{\widetilde P}}
 {(\widetilde P^2)^2}.
\end{equation*}
在领先共线幂次下写 $\mathcal S=A(z,K^2)\slashed P+\text{次领项}$。用 $\slashed n$ 投影即可确定 $A$：
\begin{equation*}
 4(P\cdot n)A
 =\operatorname{tr}(\slashed n\mathcal S).
\end{equation*}
这里的六伽马迹不能作为黑箱略过。两个物理偏振的完备关系为
\begin{equation*}
 \sum_{a=1}^{2}\varepsilon_a^{*\mu}\varepsilon_a^{\nu}
 =-\eta^{\mu\nu}
 +\frac{k^\mu n^\nu+n^\mu k^\nu}{k\cdot n},
 \qquad k\cdot n=(1-z)P\cdot n,
\end{equation*}
其中 $n^2=0$。因此先在两个偏振指标上完成收缩：
\begin{align*}
\sum_a\slashed\varepsilon_a^*\slashed p\slashed\varepsilon_a
={}&-\gamma_\mu\slashed p\gamma^\mu
 +\frac{\slashed k\slashed p\slashed n
 +\slashed n\slashed p\slashed k}{k\cdot n}\\
={}&2\slashed p
 +\frac{\slashed k\slashed p\slashed n
 +\slashed n\slashed p\slashed k}{(1-z)P\cdot n},
\end{align*}
其中用了四维 Clifford 恒等式
$\gamma_\mu\slashed p\gamma^\mu=-2\slashed p$。

于是所需迹分成三个可以逐项复核的标量：
\begin{align*}
\mathcal T_0
&\equiv2\Tr(\slashed n\slashed{\widetilde P}\slashed p\slashed{\widetilde P}),\\
\mathcal T_1
&\equiv\frac{1}{(1-z)P\cdot n}
\Tr(\slashed n\slashed{\widetilde P}\slashed k\slashed p
\slashed n\slashed{\widetilde P}),\\
\mathcal T_2
&\equiv\frac{1}{(1-z)P\cdot n}
\Tr(\slashed n\slashed{\widetilde P}\slashed n\slashed p
\slashed k\slashed{\widetilde P}).
\end{align*}
由 Sudakov 参数化可直接读出
\begin{equation*}
 n\cdot\widetilde P=P\cdot n,
 \quad n\cdot p=zP\cdot n,
 \quad n\cdot k=(1-z)P\cdot n,
 \quad \widetilde P^2\equiv Q^2=\frac{K^2}{z(1-z)},
\end{equation*}
以及
\begin{equation*}
 \widetilde P\cdot p
 =\widetilde P\cdot k
 =p\cdot k=\frac{Q^2}{2}.
\end{equation*}
先用四伽马迹
$\Tr(\slashed a\slashed b\slashed c\slashed d)
=4(a\!\cdot b\,c\!\cdot d-a\!\cdot c\,b\!\cdot d+a\!\cdot d\,b\!\cdot c)$，得到
\begin{align*}
\mathcal T_0
&=2\times4\left[
 (n\cdot\widetilde P)(p\cdot\widetilde P)
 -(n\cdot p)\widetilde P^2
 +(n\cdot\widetilde P)(\widetilde P\cdot p)
 \right]\\
&=8(P\cdot n)Q^2(1-z).
\end{align*}
六伽马迹也逐项展开。由 Clifford 反对易关系与迹循环性，任意六个四矢量满足递推式
\begin{align*}
\Tr(\slashed a\slashed b\slashed c\slashed d\slashed e\slashed f)
={}&(a\cdot b)\Tr(\slashed c\slashed d\slashed e\slashed f)
-(a\cdot c)\Tr(\slashed b\slashed d\slashed e\slashed f)\\
&+(a\cdot d)\Tr(\slashed b\slashed c\slashed e\slashed f)
-(a\cdot e)\Tr(\slashed b\slashed c\slashed d\slashed f)\\
&+(a\cdot f)\Tr(\slashed b\slashed c\slashed d\slashed e).
\end{align*}
对第一项的分子
\(\Tr(\slashed n\slashed{\widetilde P}\slashed k\slashed p
\slashed n\slashed{\widetilde P})\)，因为 $n^2=0$，递推后只剩
\begin{align*}
={}&(n\cdot\widetilde P)\Tr(\slashed k\slashed p\slashed n\slashed{\widetilde P})
-(n\cdot k)\Tr(\slashed{\widetilde P}\slashed p\slashed n\slashed{\widetilde P})\\
&+(n\cdot p)\Tr(\slashed{\widetilde P}\slashed k\slashed n\slashed{\widetilde P})
+(n\cdot\widetilde P)\Tr(\slashed{\widetilde P}\slashed k\slashed p\slashed n).
\end{align*}
四个四伽马迹依次为
\begin{align*}
\Tr(\slashed k\slashed p\slashed n\slashed{\widetilde P})&=4z(P\cdot n)Q^2,\\
\Tr(\slashed{\widetilde P}\slashed p\slashed n\slashed{\widetilde P})&=4z(P\cdot n)Q^2,\\
\Tr(\slashed{\widetilde P}\slashed k\slashed n\slashed{\widetilde P})&=4(1-z)(P\cdot n)Q^2,\\
\Tr(\slashed{\widetilde P}\slashed k\slashed p\slashed n)&=4z(P\cdot n)Q^2.
\end{align*}
把 $n\cdot k=(1-z)P\cdot n$ 与 $n\cdot p=zP\cdot n$ 代回，第二、第三项恰好抵消，故
\begin{equation*}
\Tr(\slashed n\slashed{\widetilde P}\slashed k\slashed p
\slashed n\slashed{\widetilde P})
=8z(P\cdot n)^2Q^2.
\end{equation*}
完全同样地，对
\(\Tr(\slashed n\slashed{\widetilde P}\slashed n\slashed p
\slashed k\slashed{\widetilde P})\) 递推，有
\begin{align*}
={}&(n\cdot\widetilde P)\Tr(\slashed n\slashed p\slashed k\slashed{\widetilde P})
+(n\cdot p)\Tr(\slashed{\widetilde P}\slashed n\slashed k\slashed{\widetilde P})\\
&-(n\cdot k)\Tr(\slashed{\widetilde P}\slashed n\slashed p\slashed{\widetilde P})
+(n\cdot\widetilde P)\Tr(\slashed{\widetilde P}\slashed n\slashed p\slashed k),
\end{align*}
其中四个迹分别为
\begin{equation*}
4z(P\cdot n)Q^2,
\quad4(1-z)(P\cdot n)Q^2,
\quad4z(P\cdot n)Q^2,
\quad4z(P\cdot n)Q^2.
\end{equation*}
中间两项再次抵消，所以该六伽马迹同样等于
\(8z(P\cdot n)^2Q^2\)。除以共同的
$(1-z)P\cdot n$ 后，
\begin{align*}
\mathcal T_1
&=\frac{8(P\cdot n)Q^2z}{1-z},\\
\mathcal T_2
&=\frac{8(P\cdot n)Q^2z}{1-z}.
\end{align*}
所以总迹不是直接引用的“标准结果”，而是
\begin{align*}
\mathcal T_0+\mathcal T_1+\mathcal T_2
&=8(P\cdot n)Q^2
\left[(1-z)+\frac{2z}{1-z}\right]\\
&=\frac{8(P\cdot n)Q^2}{1-z}(1+z^2)\\
&=\frac{8(P\cdot n)K^2}{z(1-z)^2}(1+z^2).
\end{align*}
这一步清楚显示了 $1+z^2$ 的来源：它是两个横向物理偏振与费米子 Clifford 代数共同留下的自旋结构，而 $1/(1-z)$ 的额外增强随后还会与共线相空间结合。

再用
\begin{equation*}
 (\widetilde P^2)^2
 =\frac{K^4}{z^2(1-z)^2},
\end{equation*}
得到
\begin{equation*}
 A(z,K^2)
 =\frac{2z(1+z^2)}{K^2}.
\end{equation*}
颜色平均给 $T^aT^a=C_F\mathbf1$，于是分裂部分对任意硬过程的平方振幅只通过普适因子
\begin{equation*}
 g_s^2C_F\frac{2z(1+z^2)}{K^2}
\end{equation*}
进入。

还需把 $(n+1)$ 体相空间在共线区因子化，而不能把相空间因子当作额外规则引用。对子系统总动量
\begin{equation*}
 Q^\mu\equiv p^\mu+k^\mu=\widetilde P^\mu,
 \qquad Q^2\equiv s_{pk}>0
\end{equation*}
插入恒等分解，Lorentz 不变相空间的递推式为
\begin{equation*}
 \dd\Phi_{n+1}
 =\dd\Phi_n(\ldots,Q)\,
 \frac{\dd s_{pk}}{2\pi}\,
 \dd\Phi_2(Q;p,k).
\end{equation*}
现在只需把二体子空间的角变量换成纵向动量分数 $z$。在 $Q$ 的瞬时静止系中，两条无质量子动量大小相同，
\begin{equation*}
 |\mathbf p_*|=|\mathbf k_*|=\frac{\sqrt{s_{pk}}}{2},
\end{equation*}
所以前文已经从 Lorentz 不变测度推得的二体相空间退化为
\begin{equation*}
 \dd\Phi_2(Q;p,k)
 =\frac{1}{32\pi^2}\,\dd\Omega_*
 =\frac{1}{32\pi^2}\,\dd\varphi_*\,\dd(\cos\theta_*).
\end{equation*}
选择辅助类光方向 $n^\mu$ 定义纵向轴。由于
\begin{equation*}
 z\equiv\frac{p\cdot n}{Q\cdot n},
\end{equation*}
在 $Q$ 静止系可取
\begin{equation*}
 n_*^\mu=\nu(1,0,0,-1),
 \qquad
 p_*^\mu=\frac{\sqrt{s_{pk}}}{2}
 (1,\sin\theta_*\cos\varphi_*,\sin\theta_*\sin\varphi_*,\cos\theta_*),
\end{equation*}
其中任意正比例因子 $\nu$ 会在比值中消去。于是
\begin{align*}
 p_*\cdot n_*
 &=\frac{\nu\sqrt{s_{pk}}}{2}(1+\cos\theta_*),\\
 Q_*\cdot n_*
 &=\nu\sqrt{s_{pk}},
\end{align*}
从而
\begin{equation*}
 z=\frac{1+\cos\theta_*}{2},
 \qquad
 \dd(\cos\theta_*)=2\,\dd z.
\end{equation*}
对未观测方位角积分 $\int_0^{2\pi}\dd\varphi_*$ 后，
\begin{align*}
 \int_0^{2\pi}\dd\varphi_*\,\dd\Phi_2
 &=\frac{1}{32\pi^2}(2\pi)(2\,\dd z)\\
 &=\frac{\dd z}{8\pi}.
\end{align*}
因此母虚度与二体子空间共同给
\begin{align*}
 \frac{\dd s_{pk}}{2\pi}\,\dd\Phi_2
 &\longrightarrow
 \frac{\dd s_{pk}\,\dd z}{16\pi^2}.
\end{align*}
最后使用已经由 Sudakov 参数化严格得到的
\begin{equation*}
 s_{pk}=\frac{K^2}{z(1-z)}.
\end{equation*}
在把 $(z,K^2)$ 选作局部相空间坐标后，固定 $z$ 取微分，
\begin{equation*}
 \left.\dd s_{pk}\right|_z
 =\frac{\dd K^2}{z(1-z)},
\end{equation*}
于是
\begin{equation*}
 \boxed{
 \dd\Phi_{\rm coll}
 =\frac{1}{16\pi^2}
 \frac{\dd z\,\dd K^2}{z(1-z)}.}
\end{equation*}
这一步也说明了 $1/[z(1-z)]$ 的来源：它不是分裂顶角的动力学奇点，而是把二体角变量与母虚度改写成 $(z,K^2)$ 时产生的 Jacobian。

因此分裂概率相对于硬过程为
\begin{align*}
 \dd P_{q\to qg}
 &=g_s^2C_F\frac{2z(1+z^2)}{K^2}
 \frac{1}{16\pi^2}
 \frac{\dd z\,\dd K^2}{z(1-z)}\\
 &=\frac{\alpha_s}{2\pi}
 C_F\frac{1+z^2}{1-z}
 \frac{\dd K^2}{K^2}\,\dd z,
\end{align*}
其中用了 $g_s^2=4\pi\alpha_s$。把 $K^2$ 重新记作正文使用的 $k_\perp^2>0$，即得\eqnrefs{eq:q-to-qg-splitting-probability-dp10}。
\end{derivation}

\begin{derivation}[双对数的三角形积分区域]
对于无质量末态，$\sum_i|\mathbf p_i|=E_{\rm cm}/c$。软胶子相对最近推力轴方向的贡献为
\begin{equation*}
 v=\frac{E_g}{E_{\rm cm}}(1-|\cos\theta|)
 \simeq\frac{E_g\theta^2}{2E_{\rm cm}}=z_gu_\theta.
\end{equation*}
从夸克分裂核的软端点 $P_{qq}\simeq2C_F/(1-z)$ 出发，使用胶子分数 $z_g=1-z$，得到\eqnrefs{eq:qcd-thrust-emission-measure}。对两条腿相加给因子二；宽角软干涉与推力轴反冲对当前双对数的系数没有额外贡献。

令 $r_z=\ln(1/z_g)$、$r_\theta=\ln(1/u_\theta)$。可分辨发射条件 $z_gu_\theta>\tau$ 变成
\begin{equation*}
 r_z\geq0,\qquad r_\theta\geq0,\qquad r_z+r_\theta<\mathcal L_T.
\end{equation*}
该区域面积为 $\mathcal L_T^2/2$。因此
\begin{equation*}
 R_T=\frac{2\alpha_sC_F}{\pi}\frac{\mathcal L_T^2}{2}
 =\frac{\alpha_sC_F}{\pi}\mathcal L_T^2.
\end{equation*}
面积中的 $1/2$ 与两条辐射腿的因子二恰好抵消；二者来源独立，遗漏任意一个都会使指数相差一倍。
\end{derivation}

\begin{derivation}[实虚抵消、指数化与多发射约束]
先引入有限的软、共线分辨下限，使全发射区域 $\Omega$ 的积分 $I_\Omega=\int_\Omega\dd\mathcal P$ 有限。在领先独立发射近似中，虚修正提供 $\ee^{-I_\Omega}$，$n$ 个相同胶子的实发射测度带 $1/n!$。若仅施加逐个发射的 $v_i<\tau$ 约束，累加任意发射数得到
\begin{align*}
 \Sigma_{\max}(\tau)
 &=\ee^{-I_\Omega}\sum_{n=0}^\infty\frac1{n!}
 \left[\int_\Omega\dd\mathcal P\,\Theta(\tau-v)\right]^n\\
 &=\exp\!\left[-\int_\Omega\dd\mathcal P\,\Theta(v-\tau)\right].
\end{align*}
移除辅助红外下限以后，只留下有限的排除区积分 $R_T$。这个实虚相消使结果不依赖“总共发射了多少个无限软胶子”。

实际推力的领先软贡献是 $\tau_T=\sum_i v_i$，不是 $\max_i v_i$。对于强有序的双对数主导区域，最大项决定总和，所以 $\Sigma_{\max}$ 给出\eqnrefs{eq:qcd-thrust-sudakov} 的精度。若要保留多个同量级发射对总和的影响，应使用拉普拉斯变量 $\nu_L$：
\begin{align*}
 \mathcal Z_T(\nu_L)
 &=\ee^{-I_\Omega}\sum_{n=0}^\infty\frac1{n!}
 \prod_{i=1}^n\int_\Omega\dd\mathcal P_i\,
 \ee^{-\nu_L\sum_i v_i}\\
 &=\exp\!\left[\int_\Omega\dd\mathcal P\,(\ee^{-\nu_Lv}-1)\right],\\
 \Sigma_T(\tau)
 &=\int_{a_{\rm inv}-\ii\infty}^{a_{\rm inv}+\ii\infty}
 \frac{\dd\nu_L}{2\pi\ii\nu_L}\,
 \ee^{\nu_L\tau}\mathcal Z_T(\nu_L),\qquad a_{\rm inv}>0.
\end{align*}
这里 $a_{\rm inv}$ 为无量纲拉普拉斯逆变换路径的正实部。$\mathcal Z_T$ 在独立软发射模型内保留了加和约束，但尚不包含 QCD 的全部相关多胶子发射。系统提高对数精度需要同时处理运行耦合、非软分裂、软相关及相应硬、喷注、软函数。
\end{derivation}

\begin{derivation}[固定阶匹配与强子化误差的尺度]
重求和适用于小 $\tau$，完整固定阶结果则保留远离两喷注极限的非奇异项。若 $\Sigma_T^{\rm FO,(n)}$ 已知到 $\alpha_s^n$，一个加性匹配定义是
\begin{equation}
 \Sigma_T^{\rm match}
 =\Sigma_T^{\rm resum}+\Sigma_T^{\rm FO,(n)}
 -\left.\Sigma_T^{\rm resum}\right|_{\alpha_s^0,\ldots,\alpha_s^n}.
 \label{eq:qcd-thrust-matching}
\end{equation}
把重求和结果展开为 $\sum_{k\geq0}\alpha_s^kr_k$，最后一项逐项去掉 $k\leq n$ 的重叠贡献，故
\begin{equation*}
 \left.\Sigma_T^{\rm match}\right|_{k\leq n}
 =\Sigma_T^{\rm FO,(n)},\qquad
 \left.\Sigma_T^{\rm match}\right|_{k>n}
 =\sum_{k>n}\alpha_s^kr_k.
\end{equation*}
这同时说明不能直接把固定阶曲线与重求和曲线相加：共有的软、共线项会被重复计数。匹配时须使用相同的截面归一化、耦合方案和尺度定义。

若非微扰软动量引起的推力位移为 $\delta\tau\sim\Lambda_E/E_{\rm cm}$，其中 $\Lambda_E$ 是强子化能标，则在宽度为 $\tau$ 的两喷注尾区，其相对影响由
\begin{equation}
 \epsilon_{\rm had}\sim\frac{\Lambda_E}{E_{\rm cm}\tau}
 \label{eq:qcd-thrust-hadronization-control}
\end{equation}
控制。即使总截面的幂修正很小，小 $\tau$ 微分分布也会放大低能效应。$\epsilon_{\rm had}\ll1$ 时可以组织矩展开；接近一时必须保留非微扰软函数的形状。这是“硬过程可微扰”与“所选观测量可微扰”之间的具体尺度区别。
\end{derivation}


